\documentclass[12pt]{amsart}
\pagestyle{plain}

\usepackage{amsmath,amssymb}
\usepackage{enumerate}
\usepackage[shortlabels]{enumitem}
\setlist[enumerate]{itemsep=6pt, topsep=6pt, parsep=0pt}
\usepackage{mathtools}
\usepackage{xcolor}
\usepackage[left=0.75in,right=1in,bottom=0.9in]{geometry}
\usepackage[T1]{fontenc}
\usepackage[parfill]{parskip}
\renewcommand{\baselinestretch}{1}

% SECTION DIVIDERS
\newcommand{\divider}{
  \vspace{-1em}
  \noindent
  \raisebox{-0.15ex}{\textbullet}%
  \hspace{0.5em}%
  \rule[0.5ex]{0.95\textwidth}{1.0pt}%
  \hspace{0.5em}%
  \raisebox{-0.15ex}{\textbullet}%
  \vspace{0em}
}

\newcommand{\norm}[1]{\left\|#1\right\|}

\title{ESE 2030 FINAL EXAM : SOLUTIONS GUIDE \\ Spring 2026}
\author{prof-g}

\begin{document}
\maketitle

\begin{center}
{\bf NOTE TO STUDENTS}
\end{center}

This solutions guide is organized by course week and topic, spanning the full
semester of material. Since there are multiple versions of this exam with
permuted problem orders \emph{and} permuted answer choices, the guide does
\emph{not} identify correct answers by letter~(A)--(E).
Instead, the correct answer is \textcolor{blue}{\textbf{highlighted in blue}}
directly within each list of choices.
Match the problem text to the one you received, then read the highlighted
choice, the explanation, and the analysis of each distractor.

Each solution includes:
\begin{itemize}
\item The correct answer (highlighted in blue in the choice list)
\item Detailed mathematical explanation
\item Analysis of every wrong choice, including partial credit designations and common misconceptions
\end{itemize}

A few synthesis problems span more than one week; these are placed under the
\emph{later} of the two weeks (the week in which the synthesis becomes
available), with cross-references noted in the explanations.

\newpage

%%%%%%%%%%%%%%%%%%%%%%%%%%%%%%%%%%%%%%%%%%%%%%%%%%%%%%%%%%%
\section*{WEEK 1: SOLVING LINEAR SYSTEMS}
%%%%%%%%%%%%%%%%%%%%%%%%%%%%%%%%%%%%%%%%%%%%%%%%%%%%%%%%%%%

%----------------------------------------------------------
\subsection*{Problem: Practical advantage of the LU factorization}
%----------------------------------------------------------

Suppose a square matrix $A$ has been factored as $A = LU$. Which statement
best describes the principal practical advantage of having this factorization
in hand?

\begin{enumerate}[(A)]
\item $A = LU$ removes the need for row pivoting.
\item $A = LU$ provides $A^{-1}$ directly as $U^{-1}L^{-1}$.
\item Solving $A\mathbf{x} = \mathbf{b}$ via $A = LU$ costs fewer arithmetic operations than Gaussian elimination on the augmented matrix $[\,A \mid \mathbf{b}\,]$.
\item \textcolor{blue}{\textbf{Once $A = LU$ is computed, additional systems $A\mathbf{x} = \mathbf{b}$ are solved by forward substitution followed by back substitution.}}
\item The diagonal of $U$ reveals whether $A$ is invertible.
\end{enumerate}

\textbf{Explanation:}
The signature use of an LU factorization is to amortize the cost of factoring
across many right-hand sides. The factorization itself costs $O(n^3)$ ---
the same as a single Gaussian elimination. The payoff arrives when one wants
to solve $A\mathbf{x} = \mathbf{b}_1, A\mathbf{x} = \mathbf{b}_2, \ldots$:
each subsequent solve is just $L\mathbf{y} = \mathbf{b}$ (forward substitution)
followed by $U\mathbf{x} = \mathbf{y}$ (back substitution), two triangular
systems at $O(n^2)$ each. This is the structural insight behind LU:
triangular systems are easy, so split the hard work of factoring from the
easy work of solving.

\textbf{Trick answer:}
``Fewer arithmetic operations than a single Gaussian elimination'' captures
the intuition that LU saves work, but misattributes the source: a single
solve via LU is the \emph{same} cost as a single Gaussian elimination.
The savings come only when the factorization is re-used across multiple
right-hand sides.

\textbf{Trick answer:}
``Diagonal of $U$ reveals invertibility'' is a true consequence --- since
$L$ has unit diagonal, $\det(A) = \det(U) = \prod u_{ii}$, so $A$ is invertible
iff $U$ has all nonzero diagonal entries --- but this is a diagnostic
byproduct, not the reason one computes the factorization.

\textbf{Trick answer:}
``$A^{-1} = U^{-1}L^{-1}$'' is a technically correct identity, but in practice
one never \emph{forms} $A^{-1}$; one solves with triangular substitutions.
The whole point of LU is to avoid computing inverses.

\textbf{Trick answer:}
``Removes the need for row pivoting'' has the relationship backwards ---
pivoting is precisely why the more general $PA = LU$ form exists. LU does
not eliminate the need for pivoting; it accommodates it.

\divider

%%%%%%%%%%%%%%%%%%%%%%%%%%%%%%%%%%%%%%%%%%%%%%%%%%%%%%%%%%%
\section*{WEEK 2: ABSTRACT VECTOR SPACES}
%%%%%%%%%%%%%%%%%%%%%%%%%%%%%%%%%%%%%%%%%%%%%%%%%%%%%%%%%%%

%----------------------------------------------------------
\subsection*{Problem: Which property is NOT a vector space axiom}
%----------------------------------------------------------

Which of the following properties is NOT one of the axioms of a vector space?

\begin{enumerate}[(A)]
\item $\mathbf{u} + \mathbf{v} = \mathbf{v} + \mathbf{u}$ for every $\mathbf{u}, \mathbf{v} \in V$
\item There exists a vector $\mathbf{0} \in V$ such that $\mathbf{v} + \mathbf{0} = \mathbf{v}$ for every $\mathbf{v} \in V$
\item $1 \cdot \mathbf{v} = \mathbf{v}$ for every $\mathbf{v} \in V$
\item $c \cdot (\mathbf{u} + \mathbf{v}) = c \cdot \mathbf{u} + c \cdot \mathbf{v}$ for every scalar $c$ and $\mathbf{u}, \mathbf{v} \in V$
\item \textcolor{blue}{\textbf{$0 \cdot \mathbf{v} = \mathbf{0}$ for every $\mathbf{v} \in V$}}
\end{enumerate}

\textbf{Explanation:}
The identity $0 \cdot \mathbf{v} = \mathbf{0}$ is a \emph{theorem}, not an
axiom. It is provable from the axioms via distributivity over scalar addition:
$$0 \cdot \mathbf{v} = (0 + 0) \cdot \mathbf{v} = 0 \cdot \mathbf{v} + 0 \cdot \mathbf{v},$$
and subtracting $0 \cdot \mathbf{v}$ from both sides gives $\mathbf{0} = 0 \cdot \mathbf{v}$.
The other four are axioms: commutativity of vector addition, existence of the
additive identity, the scalar identity axiom, and distributivity of scalar
multiplication over vector addition.

The axioms are deliberately minimal. Recognizing which properties are
primitive versus which are derivable is the point of an axiomatic system.
Note that the scalar identity $1 \cdot \mathbf{v} = \mathbf{v}$ \emph{must} be
assumed as an axiom: without it, one could define a degenerate ``scalar
multiplication'' where $c \cdot \mathbf{v} = \mathbf{0}$ for all
$c, \mathbf{v}$, and most other axioms would still hold.

\textbf{Trick answer:}
The scalar identity $1 \cdot \mathbf{v} = \mathbf{v}$ looks similar in form
to $0 \cdot \mathbf{v} = \mathbf{0}$ --- both are scalar-times-vector
identities involving a specific scalar. Students may pattern-match on shape
and pick the ``$1$ version'' as derivable from the ``$0$ version'' or vice
versa. The genuine asymmetry: the $1$-version is irreducible, while the
$0$-version follows from distributivity.

\textbf{Trick answer:}
Commutativity of vector addition is an explicit axiom in every standard
formulation, not a ``computational consequence'' of structure.

\textbf{Trick answer:}
Distributivity of scalar multiplication over vector addition is an explicit
axiom; this is distinct from distributivity over scalar addition, which is
a \emph{separate} axiom. Students sometimes conflate the two and assume one
implies the other.

\textbf{Trick answer:}
Existence of the additive identity (zero vector) is an explicit axiom in
every standard formulation. Students who think of the zero vector as
``just there'' rather than as something whose existence is asserted by an
axiom might pick this.

\divider

%----------------------------------------------------------
\subsection*{Problem: Direct sum decomposition --- equivalent characterizations}
%----------------------------------------------------------

Let $V$ be a finite-dimensional vector space, and let $U, W \subseteq V$ be
subspaces satisfying $U + W = V$. Which of the following are equivalent to
$V = U \oplus W$ under this hypothesis?

\begin{enumerate}[(A)]
\item $U \cap W = \{\mathbf{0}\}$
\item $\dim(U) + \dim(W) = \dim(V)$
\item Every vector in $V$ can be written as $\mathbf{u} + \mathbf{w}$ with $\mathbf{u} \in U$, $\mathbf{w} \in W$
\item $U$ and $W$ together contain a basis for $V$
\item \textcolor{blue}{\textbf{Both (A) and (B) (the trivial-intersection condition and the dimension-count condition)}}
\end{enumerate}

\textbf{A note on the wording:}
The phrasing ``which of the following are equivalent'' is asking which of the
\emph{listed options}, taken on its own, captures the full equivalence to
$V = U \oplus W$ under the hypothesis. Both (A) and (B) are individually
equivalent to the direct-sum condition; the option naming \emph{both of them
together} is therefore the most complete answer on the list. We graded
generously where the phrasing caused a student to pick (A) or (B) alone --- those
are each correctly identified equivalences, just incomplete relative to the
combined option.

\textbf{Explanation:}
A direct sum $V = U \oplus W$ requires \emph{both} (i) every $\mathbf{v} \in V$
decomposes as $\mathbf{u} + \mathbf{w}$ with $\mathbf{u} \in U$, $\mathbf{w} \in W$,
and (ii) the decomposition is \emph{unique}. The hypothesis $U + W = V$ already
gives existence, so uniqueness is the additional requirement.

Why uniqueness is equivalent to $U \cap W = \{\mathbf{0}\}$: suppose
$\mathbf{v} = \mathbf{u}_1 + \mathbf{w}_1 = \mathbf{u}_2 + \mathbf{w}_2$ are
two decompositions. Rearranging, $\mathbf{u}_1 - \mathbf{u}_2 = \mathbf{w}_2 - \mathbf{w}_1$.
The left side lies in $U$, the right in $W$, so this common vector lies in $U \cap W$.
If $U \cap W = \{\mathbf{0}\}$, that common vector is $\mathbf{0}$ and the two
decompositions coincide. Conversely, if the intersection contains a nonzero $\mathbf{x}$,
then $\mathbf{0} = \mathbf{0} + \mathbf{0} = \mathbf{x} + (-\mathbf{x})$ gives two
decompositions of the zero vector. So uniqueness $\Leftrightarrow$ trivial intersection.

Why the dimension count is also equivalent: take a basis $\{\mathbf{u}_1, \ldots, \mathbf{u}_p\}$
of $U$ and a basis $\{\mathbf{w}_1, \ldots, \mathbf{w}_q\}$ of $W$. Their union spans $U + W = V$,
so it contains a basis of $V$ and hence has at least $\dim(V)$ vectors --- that is,
$p + q \geq \dim(V)$. The union is itself a basis of $V$ exactly when no nontrivial
linear relation $\sum a_i \mathbf{u}_i + \sum b_j \mathbf{w}_j = \mathbf{0}$ exists,
which forces $\sum a_i \mathbf{u}_i = -\sum b_j \mathbf{w}_j \in U \cap W$. So:
the union is a basis $\Leftrightarrow$ $U \cap W = \{\mathbf{0}\}$ $\Leftrightarrow$
$\dim(U) + \dim(W) = \dim(V)$. The trivial-intersection condition and the
dimension-count condition are two faces of the same fact: they are the conditions
under which a basis of $U$ glued to a basis of $W$ is a basis of $V$.

Students often treat ``direct sum'' as a synonym for ``sum,'' missing that
the directness is precisely the uniqueness of decomposition. Recognizing that
two seemingly different conditions --- one about intersection, one about
dimensions --- describe the same structural fact is the FTLA-style pattern
of equivalent characterizations applied at the level of subspaces.

\textbf{Partial credit:}
The dimension-count characterization alone --- this captures the right
structural insight but misses that the trivial-intersection condition is
equivalent under the stated hypothesis. One ladder rung below.

\textbf{Partial credit:}
The trivial-intersection characterization alone --- captures the other
equivalent insight but misses that the dimension count is equivalent under
the stated hypothesis. Same ladder rung as the partial above.

\textbf{Trick answer:}
``$U$ and $W$ together contain a basis for $V$'' sounds structurally
relevant but is too loose --- $U \cup W$ contains a basis for $V$ whenever
$U + W = V$, regardless of directness. The trap is pattern-matching on
basis vocabulary.

\textbf{Trick answer:}
``Every vector decomposes as $\mathbf{u} + \mathbf{w}$'' is just a
restatement of $U + W = V$, which was \emph{given} as a hypothesis ---
picking this means the student did not recognize that something
\emph{additional} is needed for directness. Existence without uniqueness.

\divider

%%%%%%%%%%%%%%%%%%%%%%%%%%%%%%%%%%%%%%%%%%%%%%%%%%%%%%%%%%%
\section*{WEEK 3: LINEAR TRANSFORMATIONS}
%%%%%%%%%%%%%%%%%%%%%%%%%%%%%%%%%%%%%%%%%%%%%%%%%%%%%%%%%%%

%----------------------------------------------------------
\subsection*{Problem: Quotient space --- what is an element of $V/W$?}
%----------------------------------------------------------

Let $V$ be a finite-dimensional vector space and $W \subseteq V$ a nontrivial
subspace ($W \neq V$ and $W \neq \{0\}$).
Which of the following best describes an element of the quotient space $V/W$?

\begin{enumerate}[(A)]
\item The orthogonal projection of a vector in $V$ onto the complement of $W$
\item \textcolor{blue}{\textbf{A set of vectors in $V$ that all differ from one another by an element of $W$}}
\item A vector in $V$ paired with a vector in $W$
\item A vector in $V$ that is not in $W$
\item A vector in $V$ together with a choice of representative for $W$
\end{enumerate}

\textbf{Explanation:}
An element of $V/W$ is an \emph{equivalence class}
$$[\mathbf{v}] = \mathbf{v} + W = \{\mathbf{v} + \mathbf{w} : \mathbf{w} \in W\}.$$
The relation in play is: $\mathbf{u} \sim \mathbf{v}$ iff $\mathbf{u} - \mathbf{v} \in W$.
The equivalence class $[\mathbf{v}]$ is the set of \emph{all} vectors in $V$
related to $\mathbf{v}$ in this sense --- equivalently, all vectors that
differ from $\mathbf{v}$ by something in $W$. So an element of $V/W$ is
genuinely a \emph{set} of vectors in $V$, all pairwise differing by elements
of $W$ --- not a single vector, not a pair, not a projection.

This is the structural shift students most often miss. The quotient
construction does not pick out a ``preferred representative''; it identifies
all vectors that differ by something in $W$. The zero element of $V/W$ is
the equivalence class $[\mathbf{0}] = W$ itself --- every vector in $W$
represents the zero element. Geometrically, in $V = \mathbb{R}^3$ with $W$
a line through the origin, $V/W$ is a $2$-dimensional space whose
``points'' are the parallel translates $\mathbf{v} + W$ of $W$ ---
parallel lines in $\mathbb{R}^3$.

\textbf{Trick answer:}
``A vector in $V$ that is not in $W$'' captures the intuition that
``quotienting kills $W$,'' and a non-$W$ representative does specify a
nonzero equivalence class --- but conflates an equivalence class with a
single representative and misses that vectors inside $W$ also represent
an element of $V/W$, namely the zero class.

\textbf{Trick answer:}
``A vector in $V$ paired with a vector in $W$'' --- pairs sound like
quotient-related structure, perhaps confused with direct sum $V \oplus W$ or
a product construction. The quotient is a single object derived from $V$ and
$W$, not a pair.

\textbf{Trick answer:}
``A vector in $V$ together with a choice of representative for $W$'' ---
``representative'' is the right vocabulary for the wrong question: picking
a representative describes how one \emph{names} an equivalence class, not
what the class \emph{is}.

\textbf{Trick answer:}
``Orthogonal projection of a vector in $V$ onto the complement of $W$''
imports an inner-product notion into the abstract quotient construction.
Even when $V$ has an inner product, $V/W$ and $W^\perp$ are isomorphic but
not equal --- and this problem makes no inner-product assumption.

\divider

%----------------------------------------------------------
\subsection*{Problem: FTLA dimension count for a rank-3 transformation}
%----------------------------------------------------------

Let $T: V \to W$ be a rank 3 linear transformation between finite-dimensional
vector spaces, with $\dim(V) = 7$, $\dim(W) = 4$.
Which of the following must be TRUE?

\begin{enumerate}[(A)]
\item $\dim(\ker T) = 3$ and $\mathrm{coim}(T) \cong \mathrm{im}(T)$
\item $\dim(\ker T) = 4$ and $\mathrm{coim}(T) \cong \ker(T)$
\item $\dim(\mathrm{coim}\, T) = 4$ and $\mathrm{coker}(T) \cong \mathbb{R}^1$
\item \textcolor{blue}{\textbf{$\dim(\mathrm{coim}\, T) = 3$ and $\mathrm{coker}(T) \cong \mathbb{R}^1$}}
\item $\dim(\ker T) = 4$ and $\mathrm{coker}(T) \cong \ker(T)$
\end{enumerate}

\textbf{Explanation:}
Three facts are in play:
\begin{itemize}
\item Rank-nullity: $\dim(\ker T) = \dim V - \dim(\mathrm{im}\, T) = 7 - 3 = 4$.
\item FTLA punchline: $\mathrm{coim}(T) \cong \mathrm{im}(T)$, so $\dim(\mathrm{coim}\, T) = 3$.
\item Cokernel dimension: $\dim(\mathrm{coker}\, T) = \dim W - \dim(\mathrm{im}\, T) = 4 - 3 = 1$, so $\mathrm{coker}(T) \cong \mathbb{R}^1$.
\end{itemize}
Only the choice giving $\dim(\mathrm{coim}\, T) = 3$ together with
$\mathrm{coker}(T) \cong \mathbb{R}^1$ is consistent with all three.

The FTLA gives a natural isomorphism $\mathrm{coim}(T) \cong \mathrm{im}(T)$,
but \emph{not} $\mathrm{coker}(T) \cong \ker(T)$ in general --- those have
the same dimension only when $\dim V = \dim W$, and even then the isomorphism
is not natural. This distinction is the FTLA spine.

\textbf{Partial credit:}
The choice $\dim(\mathrm{coim}\, T) = 4$ with $\mathrm{coker}(T) \cong \mathbb{R}^1$
gets the cokernel correct, but confuses $\mathrm{coim}\, T = V/\ker T$ with
the wrong dimension count: the student remembered the quotient construction
but lost the FTLA equivalence $\mathrm{coim}(T) \cong \mathrm{im}(T)$.
Twin pair to the correct answer.

\textbf{Trick answer:}
$\dim \ker T = 4$ paired with $\mathrm{coim}(T) \cong \ker(T)$: correct
kernel dimension, but pairs it with the wrong FTLA isomorphism. Coim is
naturally isomorphic to \emph{image}, not kernel --- this is the FTLA
confusion that swaps the two non-trivial pieces of the four-subspace picture.

\textbf{Trick answer:}
$\dim \ker T = 3$ paired with $\mathrm{coim}(T) \cong \mathrm{im}(T)$: has
the correct FTLA isomorphism but pairs it with the wrong kernel dimension
--- confusing $\dim \ker T$ with $\mathrm{rank}(T)$, or equivalently using
$\dim W - \mathrm{rank}$ instead of $\dim V - \mathrm{rank}$.

\textbf{Trick answer:}
$\dim \ker T = 4$ paired with $\mathrm{coker}(T) \cong \ker(T)$: correct
kernel dimension, but pairs it with the false isomorphism. The classic
FTLA-misremembering trap: dimension counts can coincide for square
domain/codomain cases, but here $\dim \ker T = 4 \neq 1 = \dim \mathrm{coker}\, T$.

\divider

%----------------------------------------------------------
\subsection*{Problem: Equal-dimension transformation --- which equivalences hold?}
%----------------------------------------------------------

Let $T: V \to W$ be a linear transformation between finite-dimensional
vector spaces with $\dim(V) = \dim(W)$. Which of the following statements
must be TRUE?

\begin{enumerate}[(A)]
\item The cokernel of $T$ is a subspace of $W$
\item If $T$ is injective, then $T$ is an isomorphism
\item If $\ker(T) = \{\mathbf{0}\}$, then $\mathrm{im}(T) = W$
\item \textcolor{blue}{\textbf{Both ``injective $\Rightarrow$ isomorphism'' and ``trivial kernel $\Rightarrow$ image is $W$''}}
\item All of the above
\end{enumerate}

\textbf{Explanation:}
Under the hypothesis $\dim(V) = \dim(W)$, injectivity and surjectivity are
equivalent --- and either, alone, implies $T$ is an isomorphism. By
rank-nullity, $\dim(\mathrm{im}\, T) = \dim(V) - \dim(\ker T)$. If $T$ is
injective then $\dim(\ker T) = 0$, so $\dim(\mathrm{im}\, T) = \dim(V) = \dim(W)$,
forcing $\mathrm{im}(T) = W$ and giving surjectivity (and thus an isomorphism).
The trivial-kernel statement says exactly this --- injective implies
surjective --- using the kernel characterization of injectivity. So the two
statements are equivalent statements of the same fact, and both hold.

The equal-dimension hypothesis collapses injectivity and surjectivity into
a single condition, just as it does for square matrices (where invertibility,
injectivity, and surjectivity all coincide). Without equal dimensions, all
three properties can come apart --- the differentiation operator
$\mathcal{D}: \mathcal{P}_3 \to \mathcal{P}_2$ is surjective but not injective.
Recognizing that the two correct statements describe the same fact via the
FTLA is the conceptual payoff.

\textbf{Trick answer:}
``Cokernel is a subspace of $W$'' is a categorical confusion:
$\mathrm{coker}(T) = W / \mathrm{im}(T)$ is a \emph{quotient} of $W$, not a
subspace. The trap is the surface intuition ``cokernel lives on the
$W$-side, so it must sit inside $W$.'' This conflates the quotient
construction with subspace inclusion. Note that $\mathrm{coker}(T)$ \emph{is}
isomorphic to a subspace of $W$ when an inner product is available, but as
defined it is a quotient, not a subspace.

\textbf{Partial credit:}
``Injective $\Rightarrow$ isomorphism'' alone --- captures one ladder rung
but misses that the trivial-kernel statement is the equivalent kernel-based
statement of the same fact.

\textbf{Partial credit:}
``Trivial kernel $\Rightarrow$ full image'' alone --- captures the FTLA
dimension count but misses that the injective-implies-isomorphism statement
is the equivalent statement in injectivity language. Twin pair to the
partial above.

\textbf{Trick answer:}
``All of the above'' would require the cokernel claim to be true; it is not.

\divider

%%%%%%%%%%%%%%%%%%%%%%%%%%%%%%%%%%%%%%%%%%%%%%%%%%%%%%%%%%%
\section*{WEEK 4: BASES AND COORDINATES}
%%%%%%%%%%%%%%%%%%%%%%%%%%%%%%%%%%%%%%%%%%%%%%%%%%%%%%%%%%%

%----------------------------------------------------------
\subsection*{Problem: Matrix representations in two bases --- relationship}
%----------------------------------------------------------

Let $T: V \to V$ be a linear transformation on a finite-dimensional vector
space, and let $\mathcal{B}$ and $\mathcal{C}$ be two ordered bases of $V$.
Let $[T]_{\mathcal{B}}$ and $[T]_{\mathcal{C}}$ denote the matrices of the
same linear transformation $T$ in these two bases. Which of the following
best describes the relationship between $[T]_{\mathcal{B}}$ and $[T]_{\mathcal{C}}$?

\begin{enumerate}[(A)]
\item They are equal, because they represent the same linear transformation.
\item They are inverses of each other, because changing basis reverses coordinates.
\item They have the same row-reduced echelon form, because row reduction does not depend on basis.
\item \textcolor{blue}{\textbf{They are similar matrices.}}
\item They have the same columns, but written in different coordinate systems.
\end{enumerate}

\textbf{Explanation:}
The linear transformation $T$ is the same coordinate-free object in both
bases, but its matrix representation depends on the choice of basis. The
two matrices are related by a change of basis: if $S$ is the change-of-coordinates
matrix from $\mathcal{C}$-coordinates to $\mathcal{B}$-coordinates, then
$$[T]_{\mathcal{C}} = S^{-1}[T]_{\mathcal{B}}S.$$
Thus $[T]_{\mathcal{B}}$ and $[T]_{\mathcal{C}}$ are similar matrices.

\textbf{Trick answer:}
``Equal'' confuses the linear transformation with its coordinate
representation. The map $T$ is the same, but its matrix depends on the basis.

\textbf{Trick answer:}
``Inverses of each other'' confuses the coordinate-change matrices with the
matrices representing $T$. The change-of-basis maps in opposite directions
are inverses of each other; the two matrices representing $T$ are similar,
not inverse.

\textbf{Trick answer:}
``Same columns, different coordinates'' --- the columns of a matrix
representation are the coordinate vectors of the images of basis vectors.
Changing basis changes both the input basis vectors and the coordinate
system used to express their images, so the columns are not generally the
same.

\textbf{Trick answer:}
``Same row-reduced echelon form'' is not the correct relationship between
two coordinate representations of the same operator. Row reduction changes
a matrix by left multiplication with elementary matrices; change of basis
for an operator uses conjugation by an invertible matrix.

\divider

%----------------------------------------------------------
\subsection*{Problem: The coordinate map and basis dependence}
%----------------------------------------------------------

Let $V$ be an $n$-dimensional vector space and fix an ordered basis $\mathcal{B}$.
Consider the coordinate map
$\Phi_{\mathcal{B}} : V \to \mathbb{R}^n$ defined by
$\Phi_{\mathcal{B}}(\mathbf{v}) = [\mathbf{v}]_{\mathcal{B}}$.
Which of the following must be true?

\begin{enumerate}[(A)]
\item $\Phi_{\mathcal{B}}$ is an isomorphism only if $V$ is literally $\mathbb{R}^n$ (not merely isomorphic to it).
\item $\Phi_{\mathcal{B}}$ is a linear transformation but generally not invertible.
\item $\Phi_{\mathcal{B}}$ is an isomorphism, and any two ordered bases $\mathcal{B}, \mathcal{C}$ produce the same coordinate map.
\item \textcolor{blue}{\textbf{$\Phi_{\mathcal{B}}$ is an isomorphism, and a different ordered basis $\mathcal{C}$ gives a different coordinate map $\Phi_{\mathcal{C}}$.}}
\item $\Phi_{\mathcal{B}}$ is an isomorphism only if $\mathcal{B}$ is orthonormal.
\end{enumerate}

\textbf{Explanation:}
The coordinate map $\Phi_{\mathcal{B}}$ is linear (coordinates respect addition
and scalar multiplication), injective (only the zero vector has all-zero
coordinates), and surjective (every $(c_1, \ldots, c_n) \in \mathbb{R}^n$
is the coordinate vector of $c_1 \mathbf{b}_1 + \cdots + c_n \mathbf{b}_n$).
So $\Phi_{\mathcal{B}}$ is an isomorphism for \emph{any} ordered basis
$\mathcal{B}$ of any $n$-dimensional $V$. But the isomorphism is \emph{not
canonical} --- it depends on the basis choice. For any two distinct ordered
bases $\mathcal{B} \neq \mathcal{C}$, the coordinate maps $\Phi_{\mathcal{B}}$
and $\Phi_{\mathcal{C}}$ are different functions: they disagree on at least
the basis vectors themselves, and in general they differ by the
change-of-basis matrix $\Phi_{\mathcal{C}} = P \circ \Phi_{\mathcal{B}}$.

Every finite-dimensional vector space ``becomes'' $\mathbb{R}^n$ after a
basis choice. This is what justifies treating abstract vector spaces $V$
via numerical computations on coordinates --- but the price is that
\emph{the isomorphism depends on the basis}. Coordinate-free properties
survive any basis choice; coordinate-dependent statements (like the matrix
entries of a transformation) do not.

\textbf{Trick answer:}
``Linear but generally not invertible'' recognizes linearity but
mis-remembers invertibility --- the answer for a student who has not
internalized that a basis is, by definition, exactly the data needed to
make the coordinate map bijective.

\textbf{Trick answer:}
``Two bases produce the same coordinate map'' collapses the basis-dependence
of the coordinate map. Tempts students who remember ``coordinate maps are
isomorphisms'' but treat the isomorphism as canonical --- the same conceptual
error as treating $V \cong V^*$ via an unspecified isomorphism.

\textbf{Trick answer:}
``Only if $V$ is literally $\mathbb{R}^n$'' is the answer for the student
who thinks $\mathbb{R}^n$ is special; misses that coordinate maps
\emph{produce} the identification $V \cong \mathbb{R}^n$ rather than
requiring it.

\textbf{Trick answer:}
``Only if $\mathcal{B}$ is orthonormal'' imports an inner-product condition
that is irrelevant to the coordinate map. Surface match on ``orthonormal
bases are nicer'' --- true in inner product spaces for computation, but
the coordinate map exists and is an isomorphism without any inner product
structure.

\divider

%----------------------------------------------------------
\subsection*{Problem: Similarity --- the fundamental description}
%----------------------------------------------------------

Square matrices $A$ and $B$ of the same size satisfy $B = P^{-1} A P$ for some
invertible matrix $P$. Which statement most fundamentally describes the
relationship between $A$ and $B$?

\begin{enumerate}[(A)]
\item $A$ and $B$ have the same determinant and trace.
\item $A$ and $B$ have the same image as subspaces of $\mathbb{R}^n$.
\item $A$ and $B$ have the same eigenvalues.
\item $A$ is obtained from $B$ by a change of coordinates in the codomain only.
\item \textcolor{blue}{\textbf{$A$ and $B$ represent the same linear transformation expressed in different bases.}}
\end{enumerate}

\textbf{Explanation:}
The equation $B = P^{-1} A P$ is the change-of-basis formula for a linear
transformation $T : V \to V$. If $A = [T]_{\mathcal{B}}$ and $\mathcal{C}$
is another basis with change-of-basis matrix $P$ (so columns of $P$ express
the $\mathcal{C}$-vectors in $\mathcal{B}$-coordinates), then
$[T]_{\mathcal{C}} = P^{-1} A P = B$. The fundamental content of similarity
is therefore: \textbf{same transformation, different bases}. Every other
true statement about similar matrices --- equal eigenvalues, equal determinant,
equal trace, equal rank --- is a \emph{consequence} of this fact, because
these quantities are intrinsic to the transformation $T$ itself and do not
depend on the basis chosen to represent it.

This is the conceptual seed of diagonalization: an eigenbasis is a basis in
which the matrix representation is diagonal, because in that basis $T$ acts
as a pure stretch along each axis.

\textbf{Trick answer:}
``Same eigenvalues'' recognizes the strongest invariant of similarity ---
eigenvalues are preserved because they are intrinsic to the transformation
--- but identifies a consequence rather than the meaning.

\textbf{Trick answer:}
``Same determinant and trace'' recognizes two correct invariants of
similarity, but a strictly weaker observation than the eigenvalue answer,
since equal eigenvalues already imply equal determinant and trace.

\textbf{Trick answer:}
``Change of coordinates in the codomain only'' --- the change of basis
affects \emph{both} the input and output sides, since the same basis is
used for the domain and codomain of $T : V \to V$; this is why $P$ appears
on both sides of $A$, not on one side. A change of coordinates ``on one
side only'' would correspond to a different operation, not similarity.

\textbf{Trick answer:}
``Same image as subspaces of $\mathbb{R}^n$'' is the classic similarity
trap: image, kernel, and column/row spaces depend on the choice of basis
and are \emph{not} preserved under similarity. For example, $I$ and
$P^{-1} I P = I$ are similar with identical image, but $A$ and $P^{-1} A P$
for general $A$ have images that differ as concrete subspaces of $\mathbb{R}^n$.

\divider

%%%%%%%%%%%%%%%%%%%%%%%%%%%%%%%%%%%%%%%%%%%%%%%%%%%%%%%%%%%
\section*{WEEK 5: INNER PRODUCTS}
%%%%%%%%%%%%%%%%%%%%%%%%%%%%%%%%%%%%%%%%%%%%%%%%%%%%%%%%%%%

%----------------------------------------------------------
\subsection*{Problem: Inner products on function spaces and parity}
%----------------------------------------------------------

Consider the vector space $C^\infty([a,b])$ of smooth functions on an interval
$[a,b]$ with inner product $\langle \cdot, \cdot \rangle$. Under which of the
following inner products (and intervals) is the vector $p = \cosh x$ orthogonal
to the vector $q = \sinh x$?

\begin{enumerate}[(A)]
\item \textcolor{blue}{\textbf{$\displaystyle \langle p, q \rangle = \int_{-1}^{1} p(x)\, q(x)\, dx$, on $C^\infty([-1,1])$}}
\item $\displaystyle \langle p, q \rangle = \int_{0}^{1} x\, p(x)\, q(x)\, dx$, on $C^\infty([0,1])$
\item $\displaystyle \langle p, q \rangle = \int_{0}^{1} p(x)\, q(x)\, dx$, on $C^\infty([0,1])$
\item $\displaystyle \langle p, q \rangle = \int_{-1}^{1} e^{x}\, p(x)\, q(x)\, dx$, on $C^\infty([-1,1])$
\item None of the above
\end{enumerate}

\textbf{Explanation:}
Each option is a genuine inner product on the stated function space, and
orthogonality of the same pair $(\cosh, \sinh)$ depends on which inner
product is in use. Computing each:
\begin{itemize}
\item Symmetric interval, unweighted: $\int_{-1}^{1} \cosh(x)\sinh(x)\, dx$. The integrand is odd ($\cosh$ even, $\sinh$ odd, product odd), and the interval is symmetric about $0$, so the integral is $0$. \emph{Orthogonal.}
\item $[0,1]$, unweighted: $\int_{0}^{1} \cosh(x)\sinh(x)\, dx = \tfrac{1}{2}\sinh^2(1) \neq 0$. The interval $[0,1]$ breaks the symmetry that made the previous case vanish.
\item $[0,1]$ with weight $x$: $\int_{0}^{1} x \cosh(x)\sinh(x)\, dx > 0$, since the integrand is positive on $(0,1]$.
\item $[-1,1]$ with weight $e^x$: the interval is symmetric, but the weight is neither even nor odd, so the combined integrand $e^x \cosh(x) \sinh(x)$ is not odd. A direct computation gives $e^x \cosh(x)\sinh(x) = \tfrac{1}{2} e^x \sinh(2x)$, which evaluates to a nonzero quantity.
\end{itemize}

The conceptual point: orthogonality depends on the inner product. The same
two functions can be orthogonal under one inner product and not under another.
Symmetric integration interval together with an odd \emph{combined} integrand
is the only mechanism that can force the integral to zero for free.

\textbf{Trick answer:}
The weighted symmetric option ($e^x$ on $[-1,1]$) catches students who use
the heuristic ``odd function on symmetric interval gives zero'' without
checking whether the \emph{combined} integrand (weight times product) is
odd. Here $\cosh \sinh$ is odd and the interval is symmetric, but the
weight $e^x$ is asymmetric, so the heuristic fails.

\textbf{Trick answer:}
The unweighted $[0,1]$ option matches the form of the correct answer except
for the interval. Removing the symmetry destroys the parity argument; this
catches students who pattern-match on ``$\cosh \sinh$ is odd'' without
tracking the interval.

\textbf{Trick answer:}
The weighted $[0,1]$ option compounds two issues: the interval is asymmetric
\emph{and} the weight is positive, so no cancellation can occur.

\textbf{Trick answer:}
``None of the above'' is the over-cautious choice; tempts students who fail
to verify the parity argument on the symmetric unweighted case.

\divider

%----------------------------------------------------------
\subsection*{Problem: QR decomposition --- why is $R$ upper triangular?}
%----------------------------------------------------------

Let $A$ be an $m \times n$ matrix with linearly independent columns, and let
$A = QR$ be its QR decomposition. Why must $R$ be upper triangular rather
than lower triangular?

\begin{enumerate}[(A)]
\item It is a convention; the decomposition $A = QR$ also admits a lower-triangular $R$
\item \textcolor{blue}{\textbf{Because the $k$th column of $A$ is a linear combination of the first $k$ columns of $Q$}}
\item Because $R = Q^T A$ and the transpose of a triangular matrix is upper triangular
\item Because the columns of $Q$ are pairwise orthogonal
\item Because the $k$th column of $A$ is orthogonal to the first $k-1$ columns of $Q$
\end{enumerate}

\textbf{Explanation:}
The QR decomposition records exactly what Gram-Schmidt does to the columns
of $A$. If $\mathbf{q}_1, \ldots, \mathbf{q}_n$ are the orthonormalized
vectors and $\mathbf{a}_1, \ldots, \mathbf{a}_n$ are the columns of $A$,
then the partial-span invariance of Gram-Schmidt gives
$$\mathrm{span}\{\mathbf{a}_1, \ldots, \mathbf{a}_k\} = \mathrm{span}\{\mathbf{q}_1, \ldots, \mathbf{q}_k\} \quad \text{for every } k.$$
Hence $\mathbf{a}_k$ is a linear combination of $\mathbf{q}_1, \ldots, \mathbf{q}_k$
alone, with no contribution from $\mathbf{q}_{k+1}, \ldots, \mathbf{q}_n$.
Reading $A = QR$ column-by-column, the $k$th column of $R$ has zeros below
row $k$ --- precisely the upper-triangular structure.

The $(k,k)$-entry of $R$ is the norm of the unnormalized Gram-Schmidt output
at step $k$. The above-diagonal entries $R_{j,k}$ for $j < k$ are the
projection coefficients $\langle \mathbf{a}_k, \mathbf{q}_j \rangle$ that
Gram-Schmidt subtracts off. $R$ is the bookkeeping ledger of Gram-Schmidt;
its triangularity is the recursion's visible fingerprint.

\textbf{Trick answer:}
``Columns of $Q$ are pairwise orthogonal'' is a true statement, but explains
the structure of $Q$, not the triangularity of $R$. Surface match for the
student who reaches for any QR fact without tracking which fact answers
the question.

\textbf{Trick answer:}
``$k$th column of $A$ is orthogonal to the first $k-1$ columns of $Q$'' is
a twin pair to the correct answer --- same surface (``$\mathbf{a}_k$ relates
to the first $k-1$ columns of $Q$'') with the relationship swapped from
\emph{linear combination} to \emph{orthogonality}. In fact the $k$th column
of $A$ is generally \emph{not} orthogonal to the previous columns of $Q$
--- it has nonzero components along them, which is exactly what
Gram-Schmidt subtracts off.

\textbf{Trick answer:}
``$R = Q^T A$ and transpose of triangular is upper triangular'' --- the
formula $R = Q^T A$ is correct, but the rest of the reasoning is nonsense:
the transpose of an upper-triangular matrix is lower-triangular, not the
other way around. Invokes a transpose identity that doesn't say what the
student wants it to say.

\textbf{Trick answer:}
``Convention'' is the answer for the student who has not connected the
triangularity of $R$ to anything structural and assumes it is arbitrary.
In fact the partial-span structure of Gram-Schmidt \emph{forces} $R$
upper-triangular; a lower-triangular $R$ would correspond to running
Gram-Schmidt \emph{backwards} on the columns of $A$, which is a different
decomposition.

\divider

%----------------------------------------------------------
\subsection*{Problem: Cosine similarity --- which operations are invariances?}
%----------------------------------------------------------

For nonzero vectors $\mathbf{u}, \mathbf{v} \in \mathbb{R}^n$, define
$\mathrm{sim}(\mathbf{u}, \mathbf{v}) = \dfrac{\langle \mathbf{u}, \mathbf{v} \rangle}{\norm{\mathbf{u}}\,\norm{\mathbf{v}}}$.
Which of the following operations is guaranteed to leave
$\mathrm{sim}(\mathbf{u}, \mathbf{v})$ unchanged?

\begin{enumerate}[(A)]
\item \textcolor{blue}{\textbf{Replacing $\mathbf{u}$ with $5\mathbf{u}$}}
\item Replacing $\mathbf{u}$ with $-\mathbf{u}$
\item Replacing both $\mathbf{u}$ and $\mathbf{v}$ with $\mathbf{u} + \mathbf{c}$ and $\mathbf{v} + \mathbf{c}$ for a fixed $\mathbf{c} \in \mathbb{R}^n$
\item All of the above
\item None of the above
\end{enumerate}

\textbf{Explanation:}
Cosine similarity depends only on the \emph{directions} of the input
vectors, not their magnitudes --- but it is sensitive to the sign of each
vector and to the choice of origin.

For positive scaling: $\mathrm{sim}(c\mathbf{u}, \mathbf{v}) = \frac{c \langle \mathbf{u}, \mathbf{v}\rangle}{(c \norm{\mathbf{u}}) \norm{\mathbf{v}}} = \mathrm{sim}(\mathbf{u}, \mathbf{v})$.
The factors of $c$ cancel.

Negation flips the sign:
$\mathrm{sim}(-\mathbf{u}, \mathbf{v}) = -\mathrm{sim}(\mathbf{u}, \mathbf{v})$.
The numerator picks up a sign; the denominator does not (norm is positive).
Cosine similarity is invariant under \emph{positive} scaling, not arbitrary
scaling.

Translation by a common vector $\mathbf{c}$ generally changes the similarity.
The inner product $\langle \mathbf{u} + \mathbf{c}, \mathbf{v} + \mathbf{c} \rangle = \langle \mathbf{u}, \mathbf{v} \rangle + \langle \mathbf{u}, \mathbf{c} \rangle + \langle \mathbf{c}, \mathbf{v} \rangle + \norm{\mathbf{c}}^2$
does not simplify, and the norms in the denominator change too. This is
why PCA centers data \emph{before} computing similarities --- translation
is not an invariance of cosine similarity, only of \emph{centered} cosine
similarity (a.k.a.\ Pearson correlation).

Geometrically: cosine similarity is the cosine of the angle the two vectors
make at the origin. Sliding both vectors by $\mathbf{c}$ moves them away
from the origin, which generally changes the angle they subtend at the
origin even though their relative offset is unchanged. Negation rotates a
vector by $180^\circ$, flipping the sign of the cosine.

\textbf{Trick answer:}
Negation is the naive ``scaling invariance'' answer that doesn't track the
sign restriction. Tempting because $\norm{-\mathbf{u}} = \norm{\mathbf{u}}$
--- students stop checking after the denominator.

\textbf{Trick answer:}
Translation is the deepest trap. Students who think ``cosine similarity is
just relative geometry'' import a false translation invariance --- a mental
model that confuses cosine similarity with Pearson correlation.

\textbf{Trick answer:}
``All of the above'' combines the correct positive-scaling invariance with
two false ones; tempts students who recognize the scaling case and assume
the ``more invariances the better.''

\textbf{Trick answer:}
``None of the above'' is the over-cautious answer; tempts students who,
having spotted that negation and translation fail, conclude that nothing is
safe and miss the positive-scaling case.

\divider

%----------------------------------------------------------
\subsection*{Problem: The adjoint of a linear transformation}
%----------------------------------------------------------

Let $V$ and $W$ be inner product spaces and let $T: V \to W$ be a linear
transformation. The adjoint $T^*$ is the linear transformation characterized
by which of the following?

\begin{enumerate}[(A)]
\item $T^* : V \to W$ with $T^* \circ T = I_V$
\item $T^* : W \to V$ with $T(T^*(\mathbf{w})) = \mathbf{w}$ for all $\mathbf{w} \in W$
\item $T^* : W \to V$ with $\langle T(\mathbf{v}), T^*(\mathbf{w}) \rangle = \langle \mathbf{v}, \mathbf{w} \rangle$ for all $\mathbf{v} \in V$, $\mathbf{w} \in W$
\item \textcolor{blue}{\textbf{$T^* : W \to V$ with $\langle T(\mathbf{v}), \mathbf{w} \rangle = \langle \mathbf{v}, T^*(\mathbf{w}) \rangle$ for all $\mathbf{v} \in V$, $\mathbf{w} \in W$}}
\item $T^* : V \to W$ with $\langle T(\mathbf{v}), \mathbf{w} \rangle = \langle \mathbf{v}, T^*(\mathbf{w}) \rangle$ for all $\mathbf{v} \in V$, $\mathbf{w} \in W$
\end{enumerate}

\textbf{Explanation:}
The adjoint $T^*$ goes in the opposite direction to $T$, so $T^* : W \to V$.
Its defining property is the cross-pairing identity
$$\langle T(\mathbf{v}), \mathbf{w} \rangle_W = \langle \mathbf{v}, T^*(\mathbf{w}) \rangle_V$$
--- moving $T$ across the inner product produces $T^*$ on the other side.

Type-signature picture: $T(\mathbf{v}) \in W$ pairs with $\mathbf{w} \in W$
on the left, and $\mathbf{v} \in V$ pairs with $T^*(\mathbf{w}) \in V$ on
the right. Each inner product is taken in the space where both arguments
live.

\textbf{Trick answer:}
The same identity with $T^* : V \to W$ has the correct cross-pairing
identity but the wrong direction. Catches students who memorize the
identity but not which space $T^*$ acts between.

\textbf{Trick answer:}
$\langle T(\mathbf{v}), T^*(\mathbf{w}) \rangle = \langle \mathbf{v}, \mathbf{w} \rangle$
is the \emph{isometry} condition --- preserving inner products --- not the
adjoint condition. Surface match for students who confuse adjoint with the
inverse of an orthogonal transformation.

\textbf{Trick answer:}
$T^* \circ T = I_V$ with $T^* : V \to W$ is the left-inverse property ---
confuses adjoint with inverse, with the wrong type signature on top of it.

\textbf{Trick answer:}
$T(T^*(\mathbf{w})) = \mathbf{w}$ is the right-inverse property --- same
confusion as the inverse trap with the right type signature; tempts
students who got the direction right but think ``adjoint = inverse on the
codomain side.''

\divider

%----------------------------------------------------------
\subsection*{Problem: Gram-Schmidt --- partial-span invariance}
%----------------------------------------------------------

Let $\{\mathbf{v}_1, \mathbf{v}_2, \ldots, \mathbf{v}_n\}$ be a linearly
independent set in an inner product space $V$, and let
$\{\mathbf{u}_1, \mathbf{u}_2, \ldots, \mathbf{u}_n\}$ be the orthogonal set
produced by applying the Gram-Schmidt process (in that order). Which of the
following must be true for every $k$ with $1 \leq k \leq n$?

\begin{enumerate}[(A)]
\item $\mathbf{u}_k = \mathbf{v}_k$
\item $\norm{\mathbf{u}_k} = \norm{\mathbf{v}_k}$
\item \textcolor{blue}{\textbf{$\mathrm{span}\{\mathbf{u}_1, \ldots, \mathbf{u}_k\} = \mathrm{span}\{\mathbf{v}_1, \ldots, \mathbf{v}_k\}$}}
\item $\mathbf{u}_k$ is a positive scalar multiple of $\mathbf{v}_k$
\item $\mathbf{u}_k \perp \mathbf{v}_k$
\end{enumerate}

\textbf{Explanation:}
The defining feature of Gram-Schmidt is that it builds the orthogonal vectors
recursively by subtracting from $\mathbf{v}_k$ its projections onto
$\mathbf{u}_1, \ldots, \mathbf{u}_{k-1}$:
$$\mathbf{u}_k = \mathbf{v}_k - \sum_{j=1}^{k-1} \frac{\langle \mathbf{v}_k, \mathbf{u}_j\rangle}{\langle \mathbf{u}_j, \mathbf{u}_j\rangle}\, \mathbf{u}_j.$$
Each $\mathbf{u}_k$ lies in $\mathrm{span}\{\mathbf{v}_1, \ldots, \mathbf{v}_k\}$,
and conversely each $\mathbf{v}_k$ can be recovered from the $\mathbf{u}$'s.
So the partial spans match for every $k$. This is why Gram-Schmidt produces
a \emph{compatible} orthogonal basis --- it orthogonalizes without disturbing
the nested-subspace structure of the original list.

Geometric picture: at step $k$, Gram-Schmidt takes $\mathbf{v}_k$ and removes
its component inside the subspace already spanned by $\mathbf{v}_1, \ldots, \mathbf{v}_{k-1}$.
The leftover is $\mathbf{u}_k$, which lives in the same partial span as
$\mathbf{v}_k$ but is now orthogonal to everything that came before. This
nested-span invariance is exactly the structure that makes the QR decomposition
$A = QR$ work, with $R$ upper-triangular precisely because $\mathbf{v}_k$
involves only $\mathbf{u}_1, \ldots, \mathbf{u}_k$.

\textbf{Trick answer:}
$\mathbf{u}_k = \mathbf{v}_k$ is true only at $k=1$; in general
$\mathbf{u}_k \neq \mathbf{v}_k$ because Gram-Schmidt subtracts the
projections onto previous vectors. The naive answer for a student who
treats Gram-Schmidt as ``relabeling.''

\textbf{Trick answer:}
$\norm{\mathbf{u}_k} = \norm{\mathbf{v}_k}$ --- norms are not preserved;
subtracting projections strictly shortens the vector unless $\mathbf{v}_k$
is already orthogonal to the previous span. Surface match on the QR fact
that the diagonal of $R$ encodes these length changes.

\textbf{Trick answer:}
$\mathbf{u}_k \perp \mathbf{v}_k$ mixes up which vectors are orthogonal to
which. Gram-Schmidt makes $\mathbf{u}_k \perp \mathbf{u}_j$ for $j < k$,
not $\mathbf{u}_k \perp \mathbf{v}_k$. In fact
$\langle \mathbf{u}_k, \mathbf{v}_k \rangle = \norm{\mathbf{u}_k}^2 > 0$
since $\mathbf{v}_k = \mathbf{u}_k + (\text{stuff orthogonal to } \mathbf{u}_k)$.

\textbf{Trick answer:}
``$\mathbf{u}_k$ is a positive scalar multiple of $\mathbf{v}_k$'' is true
at $k=1$ only; ignores that for $k \geq 2$, $\mathbf{u}_k$ is $\mathbf{v}_k$
\emph{minus} other vectors, not a scalar multiple.

\divider

%%%%%%%%%%%%%%%%%%%%%%%%%%%%%%%%%%%%%%%%%%%%%%%%%%%%%%%%%%%
\section*{WEEK 6: ORTHOGONAL DECOMPOSITION}
%%%%%%%%%%%%%%%%%%%%%%%%%%%%%%%%%%%%%%%%%%%%%%%%%%%%%%%%%%%

%----------------------------------------------------------
\subsection*{Problem: The pseudoinverse --- minimum-norm least-squares}
%----------------------------------------------------------

Let $A$ be an $m \times n$ matrix, let $\mathbf{b} \in \mathbb{R}^m$, and let
$\hat{\mathbf{x}} = A^\dagger \mathbf{b}$, where $A^\dagger$ denotes the
pseudoinverse. Which of the following best characterizes $\hat{\mathbf{x}}$?

\begin{enumerate}[(A)]
\item $\hat{\mathbf{x}} = (A^T A)^{-1} A^T \mathbf{b}$.
\item $\hat{\mathbf{x}}$ is the unique minimizer of $\norm{A\mathbf{x} - \mathbf{b}}$ in $\mathrm{im}(A)^\perp$.
\item $\hat{\mathbf{x}}$ is the unique vector in $\mathbb{R}^n$ satisfying $A\hat{\mathbf{x}} = \mathbf{b}$.
\item $\hat{\mathbf{x}}$ is the unique vector in $\mathbb{R}^n$ such that $A\hat{\mathbf{x}}$ is the orthogonal projection of $\mathbf{b}$ onto $\mathrm{im}(A)$.
\item \textcolor{blue}{\textbf{$\hat{\mathbf{x}}$ is the unique minimizer of $\norm{A\mathbf{x} - \mathbf{b}}$ in $\ker(A)^\perp$.}}
\end{enumerate}

\textbf{Explanation:}
The pseudoinverse $A^\dagger$ produces, for any $\mathbf{b}$, the unique
vector $\hat{\mathbf{x}}$ that simultaneously (i) minimizes the residual
$\norm{A\mathbf{x} - \mathbf{b}}$, and (ii) has minimum norm among all such
minimizers. When $A$ has a nontrivial kernel, the set of residual minimizers
is the affine subspace $\hat{\mathbf{x}}_0 + \ker(A)$; the minimum-norm
choice is the unique representative in $\ker(A)^\perp$. Equivalently,
$\hat{\mathbf{x}}$ is the unique minimizer of $\norm{A\mathbf{x} - \mathbf{b}}$
that lies in $\ker(A)^\perp$ --- and existence-and-uniqueness on this
subspace is the cleanest characterization, valid whether or not exact
solutions to $A\mathbf{x} = \mathbf{b}$ exist.

Geometric picture: $A$ collapses $\ker(A)$ to zero, so any component of
$\mathbf{x}$ inside $\ker(A)$ adds norm without changing $A\mathbf{x}$.
The minimum-norm solution lives where there is nothing to throw away ---
perpendicular to the kernel. This is precisely the subspace
$\mathrm{coim}(A) = \ker(A)^\perp$, on which $A$ acts as an isomorphism
onto $\mathrm{im}(A)$.

\textbf{Trick answer:}
``Unique vector with $A\hat{\mathbf{x}} = \mathbf{b}$'' treats $A^\dagger$
as if it were $A^{-1}$. Wrong on two counts --- there may be no solution
to $A\mathbf{x} = \mathbf{b}$ at all (so the residual must be minimized
rather than zeroed), and even when solutions exist they need not be unique.

\textbf{Partial credit:}
``Unique vector such that $A\hat{\mathbf{x}}$ is the projection of
$\mathbf{b}$ onto $\mathrm{im}(A)$'' recognizes the projection geometry
--- $A\hat{\mathbf{x}}$ \emph{is} the orthogonal projection of $\mathbf{b}$
onto $\mathrm{im}(A)$ --- but misses the uniqueness issue when
$\ker(A) \neq \{\mathbf{0}\}$. Many $\mathbf{x}$ satisfy
$A\mathbf{x} = \Pi_{\mathrm{im}(A)}(\mathbf{b})$ (any two differ by an
element of $\ker(A)$), so this characterization is non-unique in general.
The student has the right geometric intuition but missed the
minimum-norm refinement that picks out a single $\hat{\mathbf{x}}$ from
the affine fiber.

\textbf{Trick answer:}
``Unique minimizer in $\mathrm{im}(A)^\perp$'' mixes domain and codomain
subspaces: $\mathrm{im}(A)^\perp \subseteq \mathbb{R}^m$, the wrong space
for $\hat{\mathbf{x}}$, which lives in $\mathbb{R}^n$. A vector in
$\mathbb{R}^m$ cannot be a minimizer over inputs $\mathbf{x}$. Surface
match on ``perp of im'' without tracking which space.

\textbf{Partial credit:}
The formula $(A^T A)^{-1} A^T \mathbf{b}$ is correct \emph{when $A$ has
full column rank}, so it gives the right answer in that special case.
Misses two structural points: the formula is conditional on full column
rank (otherwise $A^T A$ is singular), and the formula doesn't describe
what $\hat{\mathbf{x}}$ \emph{is}, only how to compute it in a special
regime.

\divider

%----------------------------------------------------------
\subsection*{Problem: Solvability of $A\mathbf{x} = \mathbf{b}$}
%----------------------------------------------------------

\textbf{A note from prof-g:} A handful of you saw a printed version of this
problem in which choice (A) read ``All three of (A), (B), (C) are equivalent
to solvability'' --- a self-referential phrasing that resulted from a bug in
the problem-shuffling script. The intent was for one option to assert that
all three of the individual conditions are equivalent to solvability. I'm
sorry for the confusion this caused. Grading on this problem will be lenient
to reflect the issue, and any partial-credit answer will be marked accordingly.
Below I describe the problem in its intended form.

Let $A \in \mathbb{R}^{m \times n}$ and $\mathbf{b} \in \mathbb{R}^m$. Which
of the following is equivalent to ``the system $A\mathbf{x} = \mathbf{b}$ has
at least one solution''?

\begin{enumerate}[(A)]
\item \textcolor{blue}{\textbf{All three of: (i) $\mathbf{b} \in \mathrm{im}(A)$, (ii) the augmented matrix $[A \mid \mathbf{b}]$ has the same rank as $A$, and (iii) $\mathbf{b} \perp \ker(A^T)$ are equivalent to solvability}}
\item The augmented matrix $[\,A \mid \mathbf{b}\,]$ has the same rank as $A$
\item $\mathbf{b} \in \mathrm{im}(A)$
\item $\mathbf{b} \perp \ker(A^T)$
\item Only $\mathbf{b} \in \mathrm{im}(A)$ is equivalent to solvability; the rank condition is a separate, stronger condition
\end{enumerate}

\textbf{Explanation:}
Three faces of the same fact:
\begin{itemize}
\item \emph{Definition.} $A\mathbf{x} = \mathbf{b}$ has a solution iff $\mathbf{b}$ lies in the image of $A$ --- this is just what ``has a solution'' means: $\mathbf{b}$ is some output of $A$.
\item \emph{Orthogonality.} $\mathbf{b} \in \mathrm{im}(A)$ iff $\mathbf{b} \perp \ker(A^T)$. Why: $\mathrm{im}(A)$ is spanned by the columns of $A$, and $\mathbf{b} \perp \ker(A^T)$ says $\mathbf{b}$ is orthogonal to every vector $\mathbf{y}$ satisfying $A^T \mathbf{y} = \mathbf{0}$ --- which is the same as saying $\mathbf{b}$ is orthogonal to every vector orthogonal to all columns of $A$. A vector that is orthogonal to everything orthogonal to the column span lies in the column span itself. (This is the geometric content of the cokernel structure: $\ker(A^T)$ is the codomain-side natural complement of $\mathrm{im}(A)$, and ``$\mathbf{b}$ is in the image'' $\Leftrightarrow$ ``$\mathbf{b}$ has zero component in this complement.'')
\item \emph{Rank criterion.} Appending $\mathbf{b}$ as an extra column to $A$ leaves the rank unchanged iff $\mathbf{b}$ is already in the span of the columns of $A$, i.e., iff $\mathbf{b} \in \mathrm{im}(A)$.
\end{itemize}

So all three statements are equivalent to solvability, and ``all three''
is correct. This is the spine of the FTLA: the computational view
(rank), the algebraic view (image), and the geometric view (orthogonality)
of the same condition.

The orthogonality version is the deepest of the three but the easiest to
misread as a stronger condition. Students often think ``orthogonal to
everything in $\ker(A^T)$'' is more demanding than ``in $\mathrm{im}(A)$''
--- but the two are exactly the same condition.

\textbf{Partial credit:}
$\mathbf{b} \in \mathrm{im}(A)$ alone --- recognizes the definitional
characterization, the most basic of the three views. Two rungs below
correct.

\textbf{Partial credit:}
The rank condition alone --- recognizes the rank-based computational
characterization, the workhorse for actually checking solvability by row
reduction. Twin pair to the previous partial.

\textbf{Partial credit:}
$\mathbf{b} \perp \ker(A^T)$ alone --- recognizes the orthogonal-FTLA
characterization, the deepest of the three views. One rung below correct;
a student who picked this has identified the most informative single
condition but missed that the others are equivalent restatements.

\textbf{Trick answer:}
``Only $\mathbf{b} \in \mathrm{im}(A)$ is equivalent; the rank condition is
stronger'' treats orthogonality (or rank) conditions as ``stronger'' than
image conditions because of the universal quantifier. The orthogonal FTLA
is precisely the theorem that collapses this apparent gap.

\divider

%----------------------------------------------------------
\subsection*{Problem: $\mathbf{b} \perp \mathrm{im}(A)$ --- which equation follows?}
%----------------------------------------------------------

Let $A$ be an $m \times n$ matrix. Suppose $\mathbf{b} \in \mathbb{R}^m$ has
the property that $\mathbf{b} \perp \mathrm{im}(A)$. Which of the following
must also be true of $\mathbf{b}$?

\begin{enumerate}[(A)]
\item $\mathbf{b} \in \mathrm{im}(A^T)$
\item \textcolor{blue}{\textbf{$A^T \mathbf{b} = \mathbf{0}$}}
\item $A \mathbf{b} = \mathbf{0}$
\item $\mathbf{b} \in \ker(A)$
\item Both of: $A^T \mathbf{b} = \mathbf{0}$ and $A\mathbf{b} = \mathbf{0}$
\end{enumerate}

\textbf{Explanation:}
The condition $\mathbf{b} \perp \mathrm{im}(A)$ says $\mathbf{b}$ is orthogonal
to every vector in the image of $A$. Since $\mathrm{im}(A)$ is spanned by the
columns of $A$, this is equivalent to $\mathbf{b}$ being orthogonal to every
column of $A$:
$$\mathbf{b} \perp \mathbf{a}_j \;\Leftrightarrow\; \mathbf{a}_j^T \mathbf{b} = 0 \quad \text{for } j = 1, \ldots, n.$$
Stacking these $n$ equations gives $A^T \mathbf{b} = \mathbf{0}$. So
$\mathbf{b} \perp \mathrm{im}(A)$ if and only if $A^T \mathbf{b} = \mathbf{0}$.
This is the direct, by-definition argument --- no FTLA machinery required.

Where this fits with the FTLA picture we developed: $\mathrm{im}(A)^\perp$ is
the natural representative of the \emph{cokernel} of $A$, and the pseudoinverse
$A^\dagger$ acts on a vector $\mathbf{c} \in \mathbb{R}^m$ by first projecting
it orthogonally onto $\mathrm{im}(A)$ --- discarding precisely the component
that lies in $\mathrm{im}(A)^\perp$ --- and then pulling that projection back
to $\ker(A)^\perp$. So the vectors in $\mathrm{im}(A)^\perp$ are exactly the
vectors that the pseudoinverse \emph{annihilates} on the codomain side. The
condition $A^T \mathbf{b} = \mathbf{0}$ is just the column-orthogonality
restatement of ``$\mathbf{b}$ lies in this discarded direction.''

Geometric picture: $\mathbf{b}$ lives in the codomain $\mathbb{R}^m$. The
orthogonal complement of $\mathrm{im}(A)$ inside $\mathbb{R}^m$ is the
codomain-side natural complement --- it does not involve the domain
$\mathbb{R}^n$ at all.

\textbf{Trick answer:}
$\mathbf{b} \in \ker(A)$ --- $\ker(A)$ lives in the domain $\mathbb{R}^n$,
not in $\mathbb{R}^m$ where $\mathbf{b}$ lives. A domain/codomain confusion.

\textbf{Trick answer:}
$\mathbf{b} \in \mathrm{im}(A^T)$ --- $\mathrm{im}(A^T) \subseteq \mathbb{R}^n$,
the wrong space; this is the orthogonal complement of $\ker(A)$ in the
domain, not of $\mathrm{im}(A)$ in the codomain.

\textbf{Trick answer:}
$A\mathbf{b} = \mathbf{0}$ --- $A\mathbf{b}$ is not even defined when
$\mathbf{b} \in \mathbb{R}^m$ unless $m = n$. Surface match on
``orthogonal $\Rightarrow$ kills under multiplication'' without tracking
which matrix to multiply by.

\textbf{Trick answer:}
``Both equations'' combines a true statement ($A^T\mathbf{b} = \mathbf{0}$)
with a wrong-domain statement ($A\mathbf{b} = \mathbf{0}$); tempts students
who recognized the correct condition but also misremembered the other as
a parallel fact.

\divider

%----------------------------------------------------------
\subsection*{Problem: Orthogonal complement --- which properties hold?}
%----------------------------------------------------------

Let $V$ be a finite-dimensional inner product space and let $W \subseteq V$
be a subspace. Which of the following must be true?

\begin{enumerate}[(A)]
\item $\dim(W) + \dim(W^\perp) = \dim(V)$
\item $W \cap W^\perp = \{\mathbf{0}\}$
\item \textcolor{blue}{\textbf{All three of: $\dim(W) + \dim(W^\perp) = \dim(V)$, $\;W \cap W^\perp = \{\mathbf{0}\}$, and $(W^\perp)^\perp = W$}}
\item $(W^\perp)^\perp = W$
\item Only the dimension count and trivial-intersection conditions hold; $(W^\perp)^\perp = W$ requires additional structure
\end{enumerate}

\textbf{Explanation:}
In a finite-dimensional inner product space, all three statements hold for
every subspace $W$:
\begin{itemize}
\item The dimension count follows from the orthogonal decomposition $V = W \oplus W^\perp$.
\item Trivial intersection is immediate: any $\mathbf{v} \in W \cap W^\perp$ satisfies $\langle \mathbf{v}, \mathbf{v} \rangle = 0$, hence $\mathbf{v} = \mathbf{0}$ by positive-definiteness of the inner product.
\item Double complement follows from the dimension count applied twice: $\dim((W^\perp)^\perp) = \dim(V) - \dim(W^\perp) = \dim(W)$, and the inclusion $W \subseteq (W^\perp)^\perp$ is automatic, so equality of dimensions forces equality of subspaces.
\end{itemize}

The orthogonal complement $\perp$ is an involution on the lattice of
subspaces of a finite-dimensional inner product space. The three properties
above are not independent observations --- they are the three faces of the
orthogonal decomposition theorem.

\textbf{Partial credit:}
The dimension-count statement alone --- correct, but the headline dimension
count only --- misses the structural facts that make the complement
well-behaved.

\textbf{Partial credit:}
The trivial-intersection statement alone --- correct, but the
trivial-intersection fact alone recognizes the direct sum is internal but
not the dimension count or the involution.

\textbf{Trick answer:}
$(W^\perp)^\perp = W$ alone --- correct as a standalone statement, but a
student who picks this over the comprehensive option has missed the other
two equivalent facts. The double-complement identity is the deepest of the
three, and the pattern-matching student grabs it without checking the others.

\textbf{Trick answer:}
The ``additional structure required'' option is a twin pair to the
double-complement choice under a fabricated caveat --- tempts the student
who half-remembers that $(W^\perp)^\perp = W$ can fail in \emph{infinite}
dimensions and overgeneralizes the caveat to the finite-dimensional setting
where it does not apply.

\divider

%----------------------------------------------------------
\subsection*{Problem: Orthogonal projection --- the closest-point characterization}
%----------------------------------------------------------

Let $V$ be a finite-dimensional inner product space and let $W \subseteq V$
be a subspace. For $\mathbf{v} \in V$, let $\Pi_W(\mathbf{v})$ denote the
orthogonal projection of $\mathbf{v}$ onto $W$. Which of the following
characterizes $\Pi_W(\mathbf{v})$?

\begin{enumerate}[(A)]
\item The unique vector $\mathbf{w} \in W$ with $\norm{\mathbf{w}} = \norm{\mathbf{v}}$
\item The unique vector $\mathbf{w} \in W$ minimizing $\norm{\mathbf{w}}$
\item The unique vector $\mathbf{w} \in V$ such that $\mathbf{v} - \mathbf{w} \in W^\perp$
\item \textcolor{blue}{\textbf{The unique vector $\mathbf{w} \in W$ minimizing $\norm{\mathbf{v} - \mathbf{w}}$}}
\item The unique vector $\mathbf{w} \in W$ such that $\mathbf{v} - \mathbf{w} \in W$
\end{enumerate}

\textbf{Explanation:}
The orthogonal projection $\Pi_W(\mathbf{v})$ is the \emph{closest point in
$W$ to $\mathbf{v}$}: it is the unique $\mathbf{w} \in W$ minimizing
$\norm{\mathbf{v} - \mathbf{w}}$. Equivalently, it is the unique
$\mathbf{w} \in W$ such that the residual $\mathbf{v} - \mathbf{w}$ lies
in $W^\perp$.

Geometric picture: drop a perpendicular from $\mathbf{v}$ to $W$; the foot
of the perpendicular is $\Pi_W(\mathbf{v})$. The residual
$\mathbf{v} - \Pi_W(\mathbf{v})$ is the perpendicular itself, lying in $W^\perp$.

\textbf{Trick answer:}
``Minimizing $\norm{\mathbf{w}}$'' --- the smallest-norm vector in $W$ is
$\mathbf{0}$ for any subspace; this characterization ignores $\mathbf{v}$
entirely. Surface match on ``minimization.''

\textbf{Trick answer:}
``$\mathbf{v} - \mathbf{w} \in W$'' instead of $W^\perp$ flips the
orthogonality condition to the wrong subspace. Twin pair to the correct
answer at the subspace level.

\textbf{Trick answer:}
``$\mathbf{w} \in V$ such that $\mathbf{v} - \mathbf{w} \in W^\perp$''
gives the correct orthogonality condition, but $\mathbf{w}$ is allowed to
range over $V$ rather than constrained to $W$. Without the constraint
$\mathbf{w} \in W$ the equation $\mathbf{v} - \mathbf{w} \in W^\perp$ has
many solutions --- every element of $\Pi_W(\mathbf{v}) + W^\perp$ --- so
uniqueness fails.

\textbf{Trick answer:}
``$\norm{\mathbf{w}} = \norm{\mathbf{v}}$'' is norm-preservation: a property
of \emph{isometries}, not projections. Projections shorten vectors in
general --- surface match on ``$\mathbf{w}$ relates to $\mathbf{v}$ by some
norm equation.''

\divider

%%%%%%%%%%%%%%%%%%%%%%%%%%%%%%%%%%%%%%%%%%%%%%%%%%%%%%%%%%%
\section*{WEEK 7: DIAGONALIZATION AND DYNAMICS}
%%%%%%%%%%%%%%%%%%%%%%%%%%%%%%%%%%%%%%%%%%%%%%%%%%%%%%%%%%%

%----------------------------------------------------------
\subsection*{Problem: Eigenvalue --- equivalent characterization}
%----------------------------------------------------------

Let $A$ be an $n \times n$ matrix and $\lambda \in \mathbb{R}$ a scalar.
Which of the following is equivalent to ``$\lambda$ is an eigenvalue of $A$''?

\begin{enumerate}[(A)]
\item $\det(A) = \lambda^n$
\item $A\mathbf{v} = \lambda \mathbf{v}$ for every $\mathbf{v} \in \mathbb{R}^n$
\item \textcolor{blue}{\textbf{$A - \lambda I$ is singular}}
\item $(A - \lambda I)\mathbf{v} = \mathbf{0}$ has only the trivial solution $\mathbf{v} = \mathbf{0}$
\item None of the above
\end{enumerate}

\textbf{Explanation:}
The chain of equivalences that licenses eigenvalue computation:
$$\lambda \text{ is an eigenvalue} \;\Leftrightarrow\; \exists\, \mathbf{v} \neq \mathbf{0} \text{ with } (A - \lambda I)\mathbf{v} = \mathbf{0} \;\Leftrightarrow\; A - \lambda I \text{ has nontrivial kernel} \;\Leftrightarrow\; A - \lambda I \text{ is singular} \;\Leftrightarrow\; \det(A - \lambda I) = 0.$$
The bridge from the geometric definition (an invariant scaling direction) to
the computational recipe (root of the characteristic polynomial) runs through
the singularity of $A - \lambda I$. That is what ``$A - \lambda I$ is singular''
names directly.

Without this bridge, the characteristic polynomial feels like an unmotivated
trick. With it, eigenvalues are exactly the scalars that ``collapse''
$A - \lambda I$ --- every other property follows.

\textbf{Trick answer:}
``$(A - \lambda I)\mathbf{v} = \mathbf{0}$ has only the trivial solution''
is the exact \emph{negation} of the eigenvalue condition. ``Only the trivial
solution'' says $\ker(A - \lambda I) = \{\mathbf{0}\}$, i.e., $A - \lambda I$
is \emph{invertible} --- which is precisely the condition under which
$\lambda$ is \emph{not} an eigenvalue. Catches students who recognize that
eigenvalues involve the equation $(A - \lambda I)\mathbf{v} = \mathbf{0}$
but lose track of the nontriviality requirement on $\mathbf{v}$. The
closest miss possible: right equation, wrong quantifier.

\textbf{Trick answer:}
``$A\mathbf{v} = \lambda \mathbf{v}$ for every $\mathbf{v}$'' replaces ``some
nonzero $\mathbf{v}$'' with ``every $\mathbf{v}$.'' Universal quantification
would mean $A = \lambda I$, a vastly stronger condition than having
$\lambda$ as an eigenvalue. Catches students who don't read quantifiers
carefully.

\textbf{Trick answer:}
$\det(A) = \lambda^n$ uses a false ``determinant is multiplicative across
$\lambda I$'' rule: $\det(A - \lambda I) \neq \det(A) - \lambda^n$, and
$\det(A) = \lambda^n$ has no general meaning. A student applying a phantom
algebraic identity lands here.

\textbf{Trick answer:}
``None of the above'' rewards the second-guesser who reads the singular
condition correctly but talks themselves out of it.

\divider

%----------------------------------------------------------
\subsection*{Problem: Spectral mapping for polynomials}
%----------------------------------------------------------

Let $A \in \mathbb{R}^{n \times n}$ and let $p(t) = c_k t^k + \cdots + c_1 t + c_0$
be a polynomial. Suppose $A\mathbf{v} = \lambda \mathbf{v}$ for a nonzero
$\mathbf{v}$. Which of the following must be TRUE?

\begin{enumerate}[(A)]
\item $p(A)$ has the same eigenvalues as $A$, since $p$ does not change the spectrum.
\item $p(A)\mathbf{v} = p(\lambda)\,\mathbf{v}$, but the eigenvalues of $p(A)$ are unrelated to those of $A$.
\item Eigenvalues of $p(A)$ are $p(\lambda_i)$ only when $A$ is diagonalizable.
\item $p(A)\mathbf{v} = \lambda\, p(\mathbf{v})$, since polynomials commute with linear maps.
\item \textcolor{blue}{\textbf{$p(A)\mathbf{v} = p(\lambda)\,\mathbf{v}$, and the eigenvectors of $A$ are eigenvectors of $p(A)$.}}
\end{enumerate}

\textbf{Explanation:}
The \emph{spectral mapping theorem for polynomials}: if $A\mathbf{v} = \lambda \mathbf{v}$,
then $A^j \mathbf{v} = \lambda^j \mathbf{v}$ for every $j \geq 0$ (induction:
$A^{j+1}\mathbf{v} = A(A^j \mathbf{v}) = A(\lambda^j \mathbf{v}) = \lambda^{j+1}\mathbf{v}$),
so applying $p$ termwise gives
$$p(A)\mathbf{v} = \sum_{j=0}^{k} c_j A^j \mathbf{v} = \sum_{j=0}^{k} c_j \lambda^j \mathbf{v} = p(\lambda)\,\mathbf{v}.$$
So $\mathbf{v}$ is an eigenvector of $p(A)$ with eigenvalue $p(\lambda)$.
The eigenvectors of $A$ remain eigenvectors of $p(A)$; the eigenvalues
transform by $p$.

This single structural fact underlies an enormous amount of the course: the
matrix exponential $e^{At} = \sum t^j A^j / j!$ has the same eigenvectors as
$A$ with eigenvalues $e^{\lambda t}$; the resolvent $(zI - A)^{-1}$ has
eigenvalues $1/(z - \lambda)$; powers $A^k$ have eigenvalues $\lambda^k$.
The rule ``$A$ acts on its own eigenvectors as multiplication by $\lambda$''
propagates through every analytic function of $A$ that the course encounters.

The result holds without any diagonalizability hypothesis --- it is a
per-eigenvector statement, not a global diagonalization statement.

\textbf{Trick answer:}
``$p(A)\mathbf{v} = p(\lambda)\mathbf{v}$, but spectra are unrelated'' gets
the per-eigenvector identity correct but then claims the spectra are
unrelated. The spectra are related precisely by $p$: every eigenvalue of
$p(A)$ is of the form $p(\lambda)$ for some eigenvalue $\lambda$ of $A$.

\textbf{Trick answer:}
``$p(A)\mathbf{v} = \lambda\, p(\mathbf{v})$'' is type confusion: $p(\mathbf{v})$
is undefined, since $p$ is a polynomial in a matrix variable, not in a
vector variable. The expression $\lambda \cdot p(\mathbf{v})$ has no
meaning. Catches students who reach for ``polynomials commute'' without
checking what objects $p$ can be applied to.

\textbf{Trick answer:}
``Only when $A$ is diagonalizable'' --- thinks diagonalizability is needed
for a per-eigenvector spectral identity. The identity holds
eigenvector-by-eigenvector; whether the eigenvectors span $\mathbb{R}^n$
is irrelevant for this single-eigenvector claim.

\textbf{Trick answer:}
``$p(A)$ has the same eigenvalues as $A$'' confuses ``preserves eigenvectors''
with ``preserves eigenvalues.'' $A^2$ has the same eigenvectors as $A$ but
eigenvalues $\lambda^2$, which generically differ from $\lambda$.

\divider

%----------------------------------------------------------
\subsection*{Problem: Matrix exponential of a nilpotent matrix}
%----------------------------------------------------------

Compute $e^{-N}$ for the matrix
$N = \begin{bmatrix} 0 & 1 & 0 \\ 0 & 0 & 1 \\ 0 & 0 & 0 \end{bmatrix}.$

\bigskip
(A) $\begin{bmatrix} 1 & 1 & 1 \\ 0 & 1 & 1 \\ 0 & 0 & 1 \end{bmatrix}$
$\quad : \quad$
(B) $\begin{bmatrix} 1 & -1 & -1/2 \\ 0 & 1 & -1 \\ 0 & 0 & 1 \end{bmatrix}$
$\quad : \quad$
(C) \textcolor{blue}{\textbf{$\begin{bmatrix} 1 & -1 & 1/2 \\ 0 & 1 & -1 \\ 0 & 0 & 1 \end{bmatrix}$}}

(D) $\begin{bmatrix} 1 & -1 & 1 \\ 0 & 1 & -1 \\ 0 & 0 & 1 \end{bmatrix}$
$\quad : \quad$
(E) $\begin{bmatrix} 1 & -1 & 0 \\ 0 & 1 & -1 \\ 0 & 0 & 1 \end{bmatrix}$
\bigskip

\textbf{Explanation:}
The matrix exponential is defined by the power series
$e^A = I + A + \frac{A^2}{2!} + \frac{A^3}{3!} + \cdots$. Applied to $-N$:
$$e^{-N} = I + (-N) + \frac{(-N)^2}{2!} + \frac{(-N)^3}{3!} + \cdots.$$
Since $N$ is nilpotent with $N^3 = 0$, the series terminates after three
terms. The key observation is that $(-N)^2 = N^2$ --- squaring cancels
the sign --- so
$$e^{-N} = I - N + \frac{N^2}{2}.$$
Computing $N^2 = \begin{bmatrix} 0 & 0 & 1 \\ 0 & 0 & 0 \\ 0 & 0 & 0 \end{bmatrix}$
and assembling gives the matrix $\begin{bmatrix} 1 & -1 & 1/2 \\ 0 & 1 & -1 \\ 0 & 0 & 1 \end{bmatrix}$.

\emph{Alternative path:} since $e^N$ and $e^{-N}$ commute and their product
is $e^{N + (-N)} = e^0 = I$, we have $e^{-N} = (e^N)^{-1}$. A student who
computes $e^N$ and inverts the resulting upper-triangular matrix arrives
at the same place. This route emphasizes that $e^A$ is always invertible
--- a foundational fact about the matrix exponential.

\textbf{Partial credit:}
The matrix $\begin{bmatrix} 1 & -1 & 1 \\ 0 & 1 & -1 \\ 0 & 0 & 1 \end{bmatrix}$
forgets the $\frac{1}{2!}$ factorial coefficient on the $N^2$ term. Knows
the series terminates and the signs are right; missed the coefficient.

\textbf{Partial credit:}
The matrix $\begin{bmatrix} 1 & -1 & 0 \\ 0 & 1 & -1 \\ 0 & 0 & 1 \end{bmatrix}$
truncates the series at $I - N$, missing the contribution from $N^2$.
Diagnoses a student who computed the first two terms and didn't check
whether higher powers vanish.

\textbf{Trick answer:}
The matrix $\begin{bmatrix} 1 & 1 & 1 \\ 0 & 1 & 1 \\ 0 & 0 & 1 \end{bmatrix}$
computes $e^{N}$ instead of $e^{-N}$, \emph{and} forgets the factorial.
The student missed the negative sign in the question entirely. No partial
credit; the sign error is fundamental.

\textbf{Trick answer:}
The matrix $\begin{bmatrix} 1 & -1 & -1/2 \\ 0 & 1 & -1 \\ 0 & 0 & 1 \end{bmatrix}$
negates the sign of \emph{every} term, including the $N^2$ term, missing
that $(-N)^2 = +N^2$. A subtle but fundamental error: squaring cancels
signs. Catches students who reflexively flip every sign rather than
thinking term-by-term.

\divider

%----------------------------------------------------------
\subsection*{Problem: Solving a linear ODE system via eigendecomposition}
%----------------------------------------------------------

Let $A$ be a $2 \times 2$ matrix with eigenvalues $\lambda_1 = 2$ and
$\lambda_2 = -1$ and corresponding eigenvectors
$\mathbf{v}_1 = \begin{pmatrix} 1 \\ 1 \end{pmatrix}, \mathbf{v}_2 = \begin{pmatrix} 1 \\ -1 \end{pmatrix}$.
What is the solution $\mathbf{x}(t)$ to the initial value problem
\[
\frac{d\mathbf{x}}{dt} = A\mathbf{x}, \qquad \mathbf{x}(0) = \begin{pmatrix} 3 \\ 1 \end{pmatrix}?
\]

\begin{enumerate}[(A)]
\item \textcolor{blue}{\textbf{$2 e^{2t} \begin{pmatrix} 1 \\ 1 \end{pmatrix} + e^{-t} \begin{pmatrix} 1 \\ -1 \end{pmatrix}$}}
\item $2 e^{2t} \begin{pmatrix} 1 \\ 1 \end{pmatrix} + e^{t} \begin{pmatrix} 1 \\ -1 \end{pmatrix}$
\item $e^{2t} \begin{pmatrix} 3 \\ 1 \end{pmatrix} + e^{-t} \begin{pmatrix} 3 \\ 1 \end{pmatrix}$
\item $e^{2t} \begin{pmatrix} 1 \\ 1 \end{pmatrix} + 2 e^{-t} \begin{pmatrix} 1 \\ -1 \end{pmatrix}$
\item $3 e^{2t} \begin{pmatrix} 1 \\ 1 \end{pmatrix} + e^{-t} \begin{pmatrix} 1 \\ -1 \end{pmatrix}$
\end{enumerate}

\textbf{Explanation:}
For diagonalizable $A$, the general solution to $\dot{\mathbf{x}} = A\mathbf{x}$
is
\[
\mathbf{x}(t) = c_1 e^{\lambda_1 t}\mathbf{v}_1 + c_2 e^{\lambda_2 t}\mathbf{v}_2,
\]
where the constants $c_1, c_2$ are the coordinates of $\mathbf{x}_0$ in the
eigenbasis. Solving $\mathbf{x}(0) = c_1 \mathbf{v}_1 + c_2 \mathbf{v}_2 = \begin{pmatrix} 3 \\ 1 \end{pmatrix}$:
\[
c_1 + c_2 = 3, \qquad c_1 - c_2 = 1 \quad \Longrightarrow \quad c_1 = 2, \; c_2 = 1.
\]
Hence $\mathbf{x}(t) = 2 e^{2t}\mathbf{v}_1 + e^{-t}\mathbf{v}_2$.

Equivalently: $\mathbf{x}(t) = e^{At}\mathbf{x}_0 = V e^{\Lambda t} V^{-1}\mathbf{x}_0$.
The vector $V^{-1}\mathbf{x}_0 = (c_1, c_2)^T$ is the initial condition
expressed in the eigenbasis; $e^{\Lambda t}$ scales each coordinate by
$e^{\lambda_i t}$; $V$ converts back to standard coordinates.

\textbf{Partial credit:}
The choice $e^{2t}\mathbf{v}_1 + 2 e^{-t}\mathbf{v}_2$ uses the correct
solution form and correct eigenvalue placement, but swaps the eigencomponent
coefficients --- gets $c_1 = 1, c_2 = 2$ instead of $c_1 = 2, c_2 = 1$.
Right framework, arithmetic slip in solving the $2 \times 2$ coordinate
system.

\textbf{Trick answer:}
The choice $3 e^{2t}\mathbf{v}_1 + e^{-t}\mathbf{v}_2$ uses the entries of
$\mathbf{x}_0 = (3,1)^T$ directly as the eigencomponent coefficients
$c_1 = 3, c_2 = 1$, skipping the eigenbasis change of coordinates. This is
the most common mistake: forgetting that $c_i$ are coordinates \emph{in
the eigenbasis}, not coordinates of $\mathbf{x}_0$ in the standard basis.

\textbf{Trick answer:}
The choice $2 e^{2t}\mathbf{v}_1 + e^{t}\mathbf{v}_2$ has a sign error on
the second eigenvalue: $e^t$ instead of $e^{-t}$. Coefficients are correct
but $\lambda_2 = -1$, not $+1$.

\textbf{Trick answer:}
The choice $e^{2t}\mathbf{x}_0 + e^{-t}\mathbf{x}_0$ is a wholesale
conceptual error: treats $\mathbf{x}_0$ as if it were the eigenvector for
both modes, replacing $\mathbf{v}_1$ and $\mathbf{v}_2$ with $\mathbf{x}_0$.
Misses the fundamental role of eigenvectors as the directions of independent
exponential evolution.

\divider

%%%%%%%%%%%%%%%%%%%%%%%%%%%%%%%%%%%%%%%%%%%%%%%%%%%%%%%%%%%
\section*{WEEK 8: EIGENVALUE COMPLEXITIES}
%%%%%%%%%%%%%%%%%%%%%%%%%%%%%%%%%%%%%%%%%%%%%%%%%%%%%%%%%%%

%----------------------------------------------------------
\subsection*{Problem: QR algorithm convergence for symmetric matrices}
%----------------------------------------------------------

Let $A \in \mathbb{R}^{n \times n}$ be a real symmetric matrix with distinct
nonzero eigenvalues. The QR algorithm produces the iteration
$$A_0 = A, \qquad A_k = Q_k R_k, \qquad A_{k+1} = R_k Q_k.$$
Each iterate $A_{k+1}$ is similar to $A_k$ (and hence to $A$), so the
eigenvalues are preserved at every step. To what does the sequence $A_k$
converge as $k \to \infty$?

\begin{enumerate}[(A)]
\item $A_k$ converges to an upper-triangular matrix whose diagonal entries are the eigenvalues of $A$.
\item $A_k$ converges to a diagonal matrix whose diagonal entries are the singular values of $A$.
\item \textcolor{blue}{\textbf{$A_k$ converges to a diagonal matrix whose diagonal entries are the eigenvalues of $A$.}}
\item $A_k$ converges to $A^T A$, since the QR iteration is built from products of $Q$ and $R$.
\item $A_k$ converges to $A$ itself, since the iteration is a similarity transformation and $A$ is the unique fixed point.
\end{enumerate}

\textbf{Explanation:}
The QR algorithm produces a sequence of matrices, each similar to $A$, that
converges to a form displaying the eigenvalues of $A$ directly on the
diagonal. For a general matrix the limit is upper-triangular (real Schur
form, with the eigenvalues on the diagonal). For a symmetric matrix --- as
here --- the limit is strictly stronger: $A_k$ converges to a \emph{diagonal}
matrix whose diagonal entries are the eigenvalues of $A$ in (typically)
descending magnitude order.

Why diagonal in the symmetric case: each iterate $A_{k+1} = Q_k^T A_k Q_k$
is itself symmetric (a similarity by an orthogonal matrix preserves
symmetry), so the limit of a sequence of symmetric upper-triangular matrices
is both symmetric and upper-triangular --- forcing all off-diagonal entries
to vanish. So the symmetric case collapses the general upper-triangular
limit to a diagonal limit.

This is the conceptual purpose of the algorithm. Similarity preserves
eigenvalues at each step; the limit, by being diagonal, \emph{displays}
those eigenvalues. Together, the two halves explain how the QR algorithm
extracts eigenvalues without ever computing a characteristic polynomial.

\textbf{Partial credit:}
``Upper-triangular with eigenvalues on the diagonal'' correctly identifies
that the limit displays eigenvalues on its diagonal --- the conceptual
heart of the convergence claim --- but states ``upper-triangular'' rather
than ``diagonal.'' Upper-triangular is the right answer for general $A$,
but the symmetric structure of this problem forces the off-diagonal entries
to zero. A student picking this has the right idea about what the QR
algorithm produces but did not exploit the symmetric hypothesis.

\textbf{Trick answer:}
``Converges to $A$ itself'' imagines the iteration as a ``settling'' that
returns the original matrix. The iteration \emph{is} a similarity
transformation at each step, but it does not fix $A$ --- instead, it
transforms $A$ along a path through the orbit of similar matrices toward a
canonical representative. The fixed points of QR iteration are the matrices
that are already (block-)triangular, not $A$ itself.

\textbf{Trick answer:}
``Diagonal of singular values'' confuses singular values with eigenvalues.
The QR algorithm extracts \emph{eigenvalues}, not singular values. For
symmetric positive-definite $A$ these happen to coincide in absolute value,
but the algorithm extracts signed eigenvalues regardless of definiteness
--- so for symmetric $A$ with negative eigenvalues the diagonal limit will
contain those negative values, not their absolute values.

\textbf{Trick answer:}
``$A^T A$'' --- $Q$ and $R$ appear in the iteration, and $A^T A = R_0^T Q_0^T Q_0 R_0 = R_0^T R_0$
at the first step, so the symbols look related --- but the iteration does
not converge to $A^T A$. For symmetric $A$, $A^T A = A^2$ has eigenvalues
$\lambda_i^2$, not $\lambda_i$, so this answer is doubly wrong: wrong limit,
wrong eigenvalue magnitudes.

\divider

%----------------------------------------------------------
\subsection*{Problem: Complex eigenvalues and oscillatory solutions}
%----------------------------------------------------------

A real matrix $A$ has a complex conjugate pair of eigenvalues $\lambda = \alpha \pm i\beta$
with $\beta \neq 0$. The system $d\mathbf{x}/dt = A\mathbf{x}$ has real basis
solutions involving $e^{\alpha t} \cos(\beta t)$ and $e^{\alpha t} \sin(\beta t)$.
Which of the following is the most fundamental reason these trigonometric
functions appear?

\begin{enumerate}[(A)]
\item \textcolor{blue}{\textbf{Euler's formula $e^{i\beta t} = \cos(\beta t) + i\sin(\beta t)$ converts the complex exponential solution into trigonometric form.}}
\item Because $A$ is real, complex eigenvalues must occur in conjugate pairs, forcing the imaginary parts of solutions to cancel.
\item The $2 \times 2$ real block associated with $\alpha \pm i\beta$ has the form $\begin{bmatrix} \alpha & -\beta \\ \beta & \alpha \end{bmatrix}$, which is a scaled rotation matrix.
\item Complex numbers do not appear in physically meaningful solutions, so they must be replaced by trigonometric functions.
\item The matrix exponential's Taylor series collapses into the series for sine and cosine when applied to the Jordan block.
\end{enumerate}

\textbf{Explanation:}
The complex-valued function $e^{\lambda t}\mathbf{v} = e^{(\alpha + i\beta)t}\mathbf{v}$
is a genuine solution to $d\mathbf{x}/dt = A\mathbf{x}$. Applying Euler's
formula, $e^{(\alpha + i\beta)t} = e^{\alpha t}(\cos(\beta t) + i\sin(\beta t))$,
so the complex solution splits into a real part and an imaginary part, each
of which is itself a real solution by linearity. The trigonometric
functions appear precisely because Euler's formula \emph{is} the bridge
between complex exponentials and oscillation. Every other appearance of
$\cos(\beta t)$ and $\sin(\beta t)$ in this setting --- in the $2 \times 2$
block formula, in the matrix-exponential Taylor series, in the structure
of damped oscillators --- derives from this single identity.

The conjugate-pair structure explains why we end up with \emph{real}
solutions, but does not by itself explain why those solutions are
\emph{trigonometric}. The block-form and Taylor-series facts are correct
\emph{consequences}: they are how Euler's formula manifests in matrix
language, not the underlying reason. Pedagogically, when students see
``where does $\cos(\beta t)$ come from?'', the answer should be Euler's
formula first; the matrix-exponential derivation is a verification of what
Euler already guarantees.

\textbf{Trick answer:}
``Real $A$ forces conjugate pairs'' is a correct fact, wrong question ---
conjugate pairing explains real-valuedness, not trigonometric structure.
A classic surface-match: ``complex conjugate'' is the right vocabulary for
a different question.

\textbf{Trick answer:}
The $2 \times 2$ real-block form is a true fact about the real Jordan
block, but the rotation interpretation is itself a consequence of Euler's
formula applied to $e^{Jt}$. Catches students who pattern-match on
``rotation'' without asking where the rotation came from.

\textbf{Partial credit:}
The Taylor-series derivation is a correct verification, but it appeals to
the structural identity $N^2 = -\beta^2 I$ (for the appropriate $2\times 2$
block) which is itself the matrix shadow of $i^2 = -1$ --- i.e., it reduces
back to Euler's formula at the scalar level. A sophisticated student might
pick this; partial credit is reasonable.

\textbf{Trick answer:}
``Complex numbers do not appear in physically meaningful solutions''
confuses physical motivation with mathematical justification. Complex
numbers do appear in solutions; we extract real solutions by linearity,
not by fiat.

\divider

%----------------------------------------------------------
\subsection*{Problem: Jordan canonical form --- algebraic vs geometric multiplicity}
%----------------------------------------------------------

A $5 \times 5$ matrix $A$ has a single eigenvalue $\lambda = 3$ with algebraic
multiplicity $5$ and geometric multiplicity $2$. Which of the following
statements is most true about the Jordan canonical form of $A$?

\begin{enumerate}[(A)]
\item \textcolor{blue}{\textbf{The Jordan form has exactly $2$ Jordan blocks, but their sizes are not determined by the given information.}}
\item The Jordan form has exactly $3$ Jordan blocks.
\item The Jordan form is diagonal, since all eigenvalues equal $3$.
\item The Jordan form is uniquely determined to be one $4 \times 4$ Jordan block plus one $1 \times 1$ block.
\item None of the above.
\end{enumerate}

\textbf{Explanation:}
This problem is about the Jordan form, but everything you need to answer it
follows from what we know about eigenspaces and shared eigenvectors.

\emph{Number of blocks = geometric multiplicity.} A Jordan block at $\lambda$
has \emph{exactly one} eigenvector direction --- the first standard basis
vector $\mathbf{e}_1$ of that block. So each Jordan block in the Jordan
form contributes exactly one independent eigenvector, and across the whole
matrix the total dimension of the eigenspace at $\lambda$ equals the number
of Jordan blocks at $\lambda$. With geometric multiplicity $2$, there are
exactly $2$ Jordan blocks.

\emph{Total size of the blocks = algebraic multiplicity.} The Jordan form is
similar to $A$, and similar matrices have the same characteristic polynomial.
Reading the characteristic polynomial off the Jordan form: each Jordan block
of size $k$ at $\lambda$ contributes a factor $(t - \lambda)^k$. So the
algebraic multiplicity of $\lambda$ is the sum of the sizes of all Jordan
blocks at $\lambda$. With algebraic multiplicity $5$, the two blocks must
have sizes summing to $5$.

\emph{Why the sizes are not pinned down.} The two ways to partition $5$
into two positive parts are $5 = 4+1$ and $5 = 3+2$. Both give matrices with
the same eigenvalue, the same algebraic multiplicity, and the same number
of independent eigenvectors --- the geometric multiplicity is $2$ in either
case --- so neither piece of information distinguishes them. The two
candidate Jordan forms are
\[
J_4(3) \oplus J_1(3) \quad \text{and} \quad J_3(3) \oplus J_2(3).
\]
These matrices behave differently in higher iterations of $(A - 3I)$ ---
$(A - 3I)^2$ and $(A - 3I)^3$ have different ranks for the two cases ---
which is what an actual computation would use to tell them apart, but the
problem doesn't give us that data.

Rule of thumb: \emph{geometric multiplicity = number of blocks; algebraic
multiplicity = total size}. The partition of block sizes is a finer
invariant that knowing eigenvalues and their two multiplicities does not
determine.

\textbf{Partial credit:}
The choice ``one $4 \times 4$ block plus one $1 \times 1$ block'' correctly
identifies one of the two valid Jordan forms and the right block count,
but asserts uniqueness; this is the ``picked the largest possible block''
instinct, which is one consistent answer but not the only one.

\textbf{Trick answer:}
``$3$ Jordan blocks'' confuses the algebraic-minus-geometric gap
($5 - 2 = 3$) with the number of blocks. That gap measures \emph{how far}
the matrix is from being diagonalizable --- it is the number of independent
directions ``missing'' from the eigenspace --- but the count of Jordan
blocks is the geometric multiplicity itself, which is $2$, not $3$.

\textbf{Trick answer:}
``Diagonal'' asserts diagonalizability from a repeated eigenvalue without
checking geometric multiplicity. With geom.\ mult.\ $= 2 \neq 5 =$ alg.\ mult.,
$A$ is not diagonalizable, so $A \neq 3I_5$.

\textbf{Trick answer:}
``None of the above'' rewards the second-guesser; the correct option here
is genuinely on the list.

\divider

%%%%%%%%%%%%%%%%%%%%%%%%%%%%%%%%%%%%%%%%%%%%%%%%%%%%%%%%%%%
\section*{WEEKS 7--8 SYNTHESIS}
%%%%%%%%%%%%%%%%%%%%%%%%%%%%%%%%%%%%%%%%%%%%%%%%%%%%%%%%%%%

%----------------------------------------------------------
\subsection*{Problem: Higher-order linear ODE --- basis of solutions}
%----------------------------------------------------------

Consider the third-order linear differential equation $p(D)x = 0$, where
$p(D) = (D - 2)^2 (D + 3)$ and $D = d/dt$. Which of the following is a basis
for the solution space?

\begin{enumerate}[(A)]
\item $\{e^{2t},\ e^{2t},\ e^{-3t}\}$
\item $\{e^{2t},\ e^{-3t}\}$
\item \textcolor{blue}{\textbf{$\{e^{2t},\ t e^{2t},\ e^{-3t}\}$}}
\item $\{e^{2t},\ e^{2t} \ln t,\ e^{-3t}\}$
\item $\{e^{-2t},\ t e^{-2t},\ e^{3t}\}$
\end{enumerate}

\textbf{Explanation:}
The characteristic polynomial $p(\lambda) = (\lambda - 2)^2 (\lambda + 3)$
has roots $\lambda = 2$ (algebraic multiplicity $2$) and $\lambda = -3$
(algebraic multiplicity $1$). Each simple root contributes one basis solution
$e^{\lambda t}$; a root of multiplicity $k$ contributes $k$ basis solutions
$e^{\lambda t}, t e^{\lambda t}, \ldots, t^{k-1} e^{\lambda t}$. So the
basis is $\{e^{2t}, t e^{2t}, e^{-3t}\}$.

Recipe: read the multiplicities off the factorization, then for each root
of multiplicity $k$, generate $k$ basis solutions by multiplying $e^{\lambda t}$
by powers of $t$ from $t^0$ to $t^{k-1}$. The total count must match the
degree of $p$, which equals the dimension of the solution space.

\textbf{Partial credit:}
$\{e^{2t}, e^{-3t}\}$ correctly identifies the two distinct exponentials
but fails to produce a \emph{basis} --- only two functions, not enough to
span a $3$-dimensional solution space. Right framework, dimension count
off by one. Captures a student who knows ``each distinct root gives an
exponential'' but missed that a repeated root contributes additional
independent solutions.

\textbf{Trick answer:}
$\{e^{2t}, e^{2t}, e^{-3t}\}$ writes $e^{2t}$ twice for the repeated root,
producing a linearly \emph{dependent} set --- this set has rank $2$, not
$3$. Diagnoses a student who knows the count must equal the multiplicity
but doesn't know how independence is achieved.

\textbf{Trick answer:}
$\{e^{-2t}, t e^{-2t}, e^{3t}\}$ sign-flips both roots: reads the factor
$(D - 2)$ as eigenvalue $-2$ and $(D + 3)$ as eigenvalue $3$. The factor
$(D - \lambda)$ corresponds to root $\lambda$, not $-\lambda$.

\textbf{Trick answer:}
$\{e^{2t}, e^{2t} \ln t, e^{-3t}\}$ multiplies by $\ln t$ instead of $t$
for the repeated root --- the polynomial ladder $1, t, t^2, \ldots$ comes
from differentiation of $e^{\lambda t}$ and the structure of Jordan blocks,
not from any logarithmic source.

\divider

%----------------------------------------------------------
\subsection*{Problem: Diagonalizability --- sufficient vs necessary conditions}
%----------------------------------------------------------

Let $A$ be an $n \times n$ matrix with real entries. Which of the following
statements about $A$ is most true?

\begin{enumerate}[(A)]
\item $A$ is diagonalizable if and only if every eigenvalue of $A$ is real.
\item If $A$ is diagonalizable, then $A$ has $n$ distinct eigenvalues, but the converse does not hold.
\item \textcolor{blue}{\textbf{If $A$ has $n$ distinct real eigenvalues, then $A$ is diagonalizable, but the converse does not hold.}}
\item $A$ is diagonalizable if and only if $A$ has $n$ distinct real eigenvalues.
\item $A$ is diagonalizable if and only if $\det(A) \neq 0$.
\end{enumerate}

\textbf{Explanation:}
Having $n$ distinct eigenvalues is \emph{sufficient} for diagonalizability:
distinct eigenvalues guarantee linearly independent eigenvectors, giving a
full eigenbasis. But it is \emph{not necessary}: the identity $I_n$ has
the single repeated eigenvalue $1$ (multiplicity $n$) and is already
diagonal. More generally, $A$ is diagonalizable iff for every eigenvalue,
geometric multiplicity equals algebraic multiplicity --- i.e., the
eigenspaces together span $\mathbb{R}^n$ (or $\mathbb{C}^n$). Distinct
eigenvalues force this automatically; repeated eigenvalues may or may not.

Geometric picture: distinct eigenvalues $\Rightarrow$ each eigenspace is
$1$-dim and they pile up to $n$ independent directions. Repeated
eigenvalues only obstruct diagonalization when the corresponding eigenspace
is too small --- a Jordan block. The implication runs one way only.

\textbf{Trick answer:}
``Diagonalizable iff $n$ distinct real eigenvalues'' asserts the converse,
the classic confusion of sufficient with necessary; the canonical
counterexample is $I_n$ or any $cI$.

\textbf{Trick answer:}
``If diagonalizable, then $n$ distinct eigenvalues'' reverses the
implication; this would mean $I$ is non-diagonalizable, which is absurd
since $I$ is already diagonal.

\textbf{Trick answer:}
``Diagonalizable iff every eigenvalue real'' makes reality of eigenvalues
independent of diagonalizability: a $2 \times 2$ rotation by
$\theta \notin \pi\mathbb{Z}$ has complex eigenvalues and is diagonalizable
over $\mathbb{C}$; a defective Jordan block can have all real eigenvalues
and not be diagonalizable.

\textbf{Trick answer:}
``Diagonalizable iff $\det(A) \neq 0$'' confuses diagonalizability with
invertibility: $I$ is diagonalizable and invertible; the zero matrix is
diagonalizable but singular; a defective $2\times 2$ Jordan block at
$\lambda = 1$ is invertible but not diagonalizable.

\divider

%----------------------------------------------------------
\subsection*{Problem: Matrix exponential of a Jordan block}
%----------------------------------------------------------

Let $J = \lambda I + N$ be a $k \times k$ Jordan block with eigenvalue
$\lambda$ on the diagonal and a single superdiagonal of $1$'s, so that $N$
is the $k \times k$ nilpotent matrix with $N^k = 0$. The matrix exponential
$e^{Jt}$ controls the solution to $d\mathbf{x}/dt = J\mathbf{x}$. Which of
the following is the most accurate description of $e^{Jt}$, and why?

\begin{enumerate}[(A)]
\item \textcolor{blue}{\textbf{$e^{Jt} = e^{\lambda t}\!\left(I + tN + \tfrac{t^2}{2!}N^2 + \cdots + \tfrac{t^{k-1}}{(k-1)!}N^{k-1}\right)$, because $\lambda I$ and $N$ commute and $N$ is nilpotent of index $k$.}}
\item $e^{Jt} = e^{\lambda I t} + e^{Nt}$, since the matrix exponential is additive over a sum of commuting matrices.
\item $e^{Jt}$ is undefined whenever $J$ is non-diagonalizable, since the matrix exponential requires an eigenbasis to compute.
\item $e^{Jt} = e^{\lambda t}\!\left(I + tN + \tfrac{t^2}{2!}N^2 + \cdots\right)$ as an infinite series, since the superdiagonal contributions never terminate.
\item $e^{Jt} = e^{\lambda t} I$, since the diagonal entry $\lambda$ governs the exponential growth rate.
\end{enumerate}

\textbf{Explanation:}
The decomposition $J = \lambda I + N$ exhibits $J$ as a sum of two
\emph{commuting} matrices --- $\lambda I$ commutes with everything --- so
the matrix exponential factors as
$$e^{Jt} = e^{(\lambda I + N) t} = e^{\lambda I t}\, e^{N t} = e^{\lambda t}\, e^{Nt}.$$
The first factor is the scalar $e^{\lambda t}$ times $I$; the second factor
is the matrix exponential of the nilpotent $N$, which by definition is the
power series $e^{Nt} = \sum_{j \geq 0} \frac{t^j}{j!} N^j$. Because
$N^k = 0$, this series terminates after $k$ terms:
$$e^{Nt} = I + tN + \tfrac{t^2}{2!} N^2 + \cdots + \tfrac{t^{k-1}}{(k-1)!} N^{k-1}.$$
Multiplying by $e^{\lambda t}$ gives the correct expression.

This is the structural origin of the polynomial-times-exponential ladder
$\{e^{\lambda t}, t e^{\lambda t}, t^2 e^{\lambda t}/2!, \ldots\}$ that
appears as the basis of solutions to a higher-order linear ODE with a
repeated root. The factor $e^{\lambda t}$ comes from $\lambda I$; the
polynomial factors $1, t, t^2/2!, \ldots, t^{k-1}/(k-1)!$ are the entries
that $N$ generates as it climbs the superdiagonal. The Jordan block is the
matrix avatar of the repeated-root behavior --- Jordan structure and the
polynomial ladder are the same phenomenon viewed from two angles.

Conceptual through-line: matrix exponentials of diagonalizable matrices
involve only $e^{\lambda t}$ scalings on each eigenvector. The \emph{new}
content here is what happens when there are not enough eigenvectors --- the
missing directions are filled in by the nilpotent $N$, and matrix
exponentiation reaches into those directions through the polynomial
$t^j N^j / j!$ ladder.

\textbf{Trick answer:}
$e^{Jt} = e^{\lambda t} I$ takes the diagonal-case answer and applies it to
the non-diagonalizable case. This would be correct if $N = 0$, but the
whole point of a Jordan block is that $N \neq 0$; ignoring $N$ ignores the
entire Jordan-chain structure.

\textbf{Trick answer:}
$e^{Jt} = e^{\lambda I t} + e^{Nt}$ invokes a false ``additivity'' rule.
The matrix exponential is multiplicative over commuting summands,
$e^{A+B} = e^A e^B$ when $AB = BA$ --- never additive. The correct
decomposition does start from $J = \lambda I + N$ with commutativity, but
produces a \emph{product} $e^{\lambda I t} e^{Nt}$, not a sum. The student
who picks this has the right ingredients but the wrong combinator.

\textbf{Partial credit:}
The infinite-series version gets the structural picture right ---
$e^{Jt} = e^{\lambda t} e^{Nt}$ with $e^{Nt}$ as a series in $N$ --- but
fails to register that $N^k = 0$ truncates the series. The student
understands the decomposition but has not exploited nilpotency. One rung
below correct: same framework, missing the finite-series payoff.

\textbf{Trick answer:}
``$e^{Jt}$ is undefined for non-diagonalizable $J$'' confuses ``computation
by diagonalization'' with the existence of $e^{Jt}$. The matrix exponential
is defined for \emph{any} square matrix via its power series; diagonalization
is a convenient shortcut when available, not a precondition.

\divider

%%%%%%%%%%%%%%%%%%%%%%%%%%%%%%%%%%%%%%%%%%%%%%%%%%%%%%%%%%%
\section*{WEEK 9: LINEAR ITERATIVE SYSTEMS}
%%%%%%%%%%%%%%%%%%%%%%%%%%%%%%%%%%%%%%%%%%%%%%%%%%%%%%%%%%%

%----------------------------------------------------------
\subsection*{Problem: Spectral Theorem for symmetric matrices}
%----------------------------------------------------------

\textbf{A note from prof-g:} A handful of you saw a printed version of this
problem in which choice (B) read ``All three of (A), (B), and (C) are
guaranteed, and each can fail for some non-symmetric real matrix'' --- a
self-referential phrasing that resulted from a bug in the problem-shuffling
script. The intent was for one option to assert that all three of the
individually-listed Spectral Theorem guarantees hold simultaneously. I'm
sorry for the confusion this caused. Grading on this problem will be lenient
to reflect the issue, and any partial-credit answer will be marked accordingly.
Below I describe the problem in its intended form.

Let $A \in \mathbb{R}^{n \times n}$ be a real symmetric matrix. The Spectral
Theorem asserts a package of structural facts about $A$ that can each fail
for a general (non-symmetric) real $n \times n$ matrix. Which of the
following is GUARANTEED by the Spectral Theorem for symmetric $A$?

\begin{enumerate}[(A)]
\item For every eigenvalue $\lambda$ of $A$, the geometric multiplicity of $\lambda$ equals its algebraic multiplicity.
\item \textcolor{blue}{\textbf{All three of: real eigenvalues, full geometric multiplicity, and an orthonormal eigenbasis --- and each can fail for some non-symmetric real matrix.}}
\item There exists an orthonormal basis of $\mathbb{R}^n$ consisting of eigenvectors of $A$.
\item Every eigenvalue of $A$ is real.
\item Only real eigenvalues and full geometric multiplicity are guaranteed; an orthonormal eigenbasis can fail when $A$ has repeated eigenvalues.
\end{enumerate}

\textbf{Explanation:}
The Spectral Theorem for real symmetric matrices asserts that
$A = Q\Lambda Q^T$ where $Q$ is orthogonal (real) and $\Lambda$ is diagonal
(real). Reading off the three guarantees:
\begin{itemize}
\item \emph{Real eigenvalues}: the diagonal entries of $\Lambda$ are real.
\item \emph{Full geometric multiplicity}: the columns of $Q$ form a basis of $\mathbb{R}^n$, so the eigenspaces collectively span $\mathbb{R}^n$ --- equivalently, geometric multiplicity equals algebraic multiplicity for every eigenvalue (no Jordan obstruction).
\item \emph{Orthonormal eigenbasis}: the columns of $Q$ are not just an eigenbasis but an orthonormal one.
\end{itemize}

Each fails for some general real matrix: real eigenvalues fail for the
rotation $\begin{pmatrix} 0 & -1 \\ 1 & 0 \end{pmatrix}$, with eigenvalues
$\pm i$; full geometric multiplicity fails for the Jordan block
$\begin{pmatrix} 2 & 1 \\ 0 & 2 \end{pmatrix}$, with algebraic multiplicity
$2$ and geometric multiplicity $1$ at $\lambda = 2$; an orthonormal
eigenbasis fails for $\begin{pmatrix} 1 & 1 \\ 0 & 2 \end{pmatrix}$, which
has an eigenbasis $\{(1,0)^T, (1,1)^T\}$ --- so it is diagonalizable ---
but no orthonormal eigenbasis. All three statements are independently
nontrivial guarantees, all three are simultaneously delivered by the
Spectral Theorem, and the comprehensive option is the only choice that
captures the full content along with the contrast.

The orthonormal-eigenbasis fact alone logically implies the other two ---
if an orthonormal eigenbasis exists, then $A = Q\Lambda Q^T$ with real $Q$
forces real $\Lambda$, and the existence of an eigenbasis is exactly full
geometric multiplicity. So the orthonormal-eigenbasis fact is the
strongest single statement on the list. The reason the comprehensive
answer is still correct is that it names the contrast with general matrices
--- the structural payoff of the theorem --- whereas the orthonormal-eigenbasis
fact alone reads like a fact about a particular $A$. The Spectral Theorem's
content is that \emph{symmetry forces this package}.

\textbf{Trick answer:}
``Orthonormal eigenbasis exists'' alone is the strongest single guarantee
on the list, logically implying the other two; a student who picked it has
identified the most informative individual fact but missed the synthesis
and the contrast with general matrices.

\textbf{Trick answer:}
``Every eigenvalue is real'' alone captures a genuine guarantee of the
theorem, but the weakest of the three on its own. Real eigenvalues alone
do not yield diagonalizability, let alone orthogonal diagonalizability.

\textbf{Trick answer:}
``Full geometric multiplicity'' alone is equivalent to diagonalizability,
a genuine guarantee, but does not address reality of the eigenvalues or
orthogonality of the eigenbasis.

\textbf{Trick answer:}
``Only real eigenvalues and full geometric multiplicity are guaranteed;
orthonormal eigenbasis can fail'' imports the Jordan-block intuition ---
``repeated eigenvalues can cause geometric multiplicity to fall short''
--- and applies it to symmetric matrices, where the Spectral Theorem
precisely \emph{prevents} this. A student who picked this has correctly
diagnosed the general behavior of Jordan blocks but failed to register
that symmetry overrides it.

\divider

%----------------------------------------------------------
\subsection*{Problem: Perron-Frobenius theorem --- complete statement}
%----------------------------------------------------------

Suppose $A$ is a square matrix with all entries strictly positive. Which
statement most completely captures what the Perron-Frobenius theorem
guarantees about the eigenstructure of $A$?

\begin{enumerate}[(A)]
\item The spectral radius $\rho(A)$ is itself an eigenvalue of $A$.
\item \textcolor{blue}{\textbf{There is a real positive eigenvalue, strictly greater in magnitude than every other eigenvalue, with a corresponding eigenvector having all strictly positive entries.}}
\item The largest magnitude eigenvalue is dominant.
\item There exists an eigenvector of $A$ with all strictly positive entries.
\item All eigenvalues of $A$ are real and positive.
\end{enumerate}

\textbf{Explanation:}
The Perron-Frobenius theorem (in its strictly-positive form) asserts a
\textbf{bundle} of facts about the dominant eigenpair of a positive matrix:
there is a single eigenvalue $\lambda_*$ that is (i) real and positive,
(ii) equal to the spectral radius $\rho(A)$, (iii) \emph{strictly} dominant
(no other eigenvalue has the same magnitude), and (iv) associated with an
eigenvector $\mathbf{v}_*$ whose entries are all strictly positive. The
correct choice packages all of these elements --- positive real, strictly
dominant, positive eigenvector --- in a single statement.

The positive eigenvector is what makes Perron-Frobenius indispensable in
ranking applications: the entries of $\mathbf{v}_*$ can be interpreted as
probabilities, populations, or scores.

\textbf{Trick answer:}
``$\rho(A)$ is itself an eigenvalue'' captures one consequence of
Perron-Frobenius --- that $\rho(A)$ is realized as an actual eigenvalue
rather than just a bound --- but says nothing about strict dominance or
the positive eigenvector. One rung of the ladder.

\textbf{Trick answer:}
``Eigenvector with positive entries'' captures the positive-eigenvector
half of the theorem but is silent about which eigenvalue this eigenvector
corresponds to. A second rung of the ladder.

\textbf{Trick answer:}
``Largest magnitude eigenvalue is dominant'' captures strict dominance ---
the gap between $\rho(A)$ and the next-largest magnitude --- but says
nothing about the sign or reality of the dominant eigenvalue, nor about
the positive eigenvector. A third rung of the ladder.

\textbf{Trick answer:}
``All eigenvalues real and positive'' over-generalizes ``positive matrix''
to ``all eigenvalues positive.'' Perron-Frobenius guarantees only that the
\emph{dominant} eigenvalue is real and positive; the other eigenvalues
can be complex or negative, provided their magnitudes are strictly smaller.
For example, the $3 \times 3$ positive matrix
$\begin{bmatrix} 1 & 1 & 1 \\ 2 & 1 & 1 \\ 1 & 2 & 1 \end{bmatrix}$ has
one dominant real eigenvalue near $3.4$ and a complex conjugate pair with
magnitude near $0.6$.

\divider

%----------------------------------------------------------
\subsection*{Problem: Stationary distribution of a positive Markov chain}
%----------------------------------------------------------

Let $P$ be a positive column-stochastic $n \times n$ matrix governing a
Markov chain
\[
\mathbf{x}_{k+1} = P \mathbf{x}_k
\]
on probability vectors in $\mathbb{R}^n$. A \emph{stationary distribution}
for $P$ is a probability vector $\boldsymbol{\pi}$ satisfying
$P\boldsymbol{\pi} = \boldsymbol{\pi}$. Which of the following is the best
Perron-Frobenius description of the stationary distribution?

\begin{enumerate}[(A)]
\item A nonzero vector in $\ker(P-I)$.
\item The limit $\lim_{k\to\infty} P^k\mathbf{x}_0$ for any probability vector $\mathbf{x}_0$.
\item A probability vector $\boldsymbol{\pi}$ such that $\boldsymbol{\pi}^T P = \boldsymbol{\pi}^T$.
\item An eigenvector of $P$ whose entries sum to $1$.
\item \textcolor{blue}{\textbf{The probability eigenvector of $P$ corresponding to the dominant eigenvalue $\lambda = 1$.}}
\end{enumerate}

\textbf{Explanation:}
Since $P$ is positive and column-stochastic, Perron-Frobenius applies cleanly.
The eigenvalue $\lambda = 1$ is the dominant eigenvalue, and its eigenspace
contains a positive eigenvector. Normalizing that positive eigenvector so
that its entries sum to $1$ gives the stationary distribution.

The correct choice captures both essential features: the vector is fixed by
the dynamics, $P\boldsymbol{\pi} = \boldsymbol{\pi}$, and it is a probability
vector. The word ``dominant'' identifies the Perron-Frobenius role of the
eigenvalue $\lambda = 1$.

\textbf{Partial credit:}
``A nonzero vector in $\ker(P-I)$'' recognizes the fixed-point equation
$P\boldsymbol{\pi} = \boldsymbol{\pi}$, equivalently
$\boldsymbol{\pi} \in \ker(P-I)$, but misses the probability normalization
and nonnegativity requirement.

\textbf{Trick answer:}
``$\boldsymbol{\pi}^T P = \boldsymbol{\pi}^T$'' uses the row-vector
convention. Under the column-stochastic convention used here, stationary
distributions are right eigenvectors of $P$, not left eigenvectors.

\textbf{Partial credit:}
``$\lim P^k \mathbf{x}_0$'' is true under the positivity hypothesis, but
describes the long-term limiting behavior rather than the Perron-Frobenius
eigenvector description of what the stationary distribution \emph{is}.

\textbf{Trick answer:}
``Eigenvector whose entries sum to $1$'' requires entries summing to $1$,
but does not specify the eigenvalue $\lambda = 1$ and does not require
nonnegative entries. So it need not describe a probability vector fixed
by the dynamics.

\divider

%%%%%%%%%%%%%%%%%%%%%%%%%%%%%%%%%%%%%%%%%%%%%%%%%%%%%%%%%%%
\section*{WEEK 10: SINGULAR VALUE DECOMPOSITION}
%%%%%%%%%%%%%%%%%%%%%%%%%%%%%%%%%%%%%%%%%%%%%%%%%%%%%%%%%%%

%----------------------------------------------------------
\subsection*{Problem: Maximum stretch of a unit vector under the SVD}
%----------------------------------------------------------

A matrix $A \in \mathbb{R}^{4 \times 3}$ has singular values $\sigma_1 = 6$,
$\sigma_2 = 3$, $\sigma_3 = 2$. Over all unit vectors
$\mathbf{x} \in \mathbb{R}^3$, what is the largest possible value of
$\norm{A\mathbf{x}}$?

\begin{enumerate}[(A)]
\item $\sqrt{6^2 + 3^2 + 2^2} = 7$
\item $36$
\item $11$
\item \textcolor{blue}{\textbf{$6$}}
\item Cannot be determined from the information given.
\end{enumerate}

\textbf{Explanation:}
The SVD gives a complete geometric picture of how $A$ acts on the unit sphere.
The right singular vectors $\{\mathbf{v}_1, \mathbf{v}_2, \mathbf{v}_3\}$ are
an orthonormal basis of $\mathbb{R}^3$, and $A$ sends each $\mathbf{v}_i$
to $\sigma_i \mathbf{u}_i$, where the $\mathbf{u}_i$ are an orthonormal set
in $\mathbb{R}^4$. So $A$ takes the unit sphere in $\mathbb{R}^3$ to an
ellipsoid in $\mathbb{R}^4$ whose three semi-axes have lengths
$\sigma_1 = 6$, $\sigma_2 = 3$, $\sigma_3 = 2$, oriented along the
$\mathbf{u}_i$.

The largest possible value of $\norm{A\mathbf{x}}$ over the unit sphere is
the longest semi-axis of this image ellipsoid. The longest semi-axis has
length $\sigma_1 = 6$, attained when $\mathbf{x} = \mathbf{v}_1$ (since
$A\mathbf{v}_1 = \sigma_1 \mathbf{u}_1$, which has norm $\sigma_1$). So
the answer is $6$.

To see directly that no other unit input does better: write any unit vector
in the right-singular-vector basis as $\mathbf{x} = c_1 \mathbf{v}_1 + c_2 \mathbf{v}_2 + c_3 \mathbf{v}_3$
with $c_1^2 + c_2^2 + c_3^2 = 1$. Then
$A\mathbf{x} = c_1 \sigma_1 \mathbf{u}_1 + c_2 \sigma_2 \mathbf{u}_2 + c_3 \sigma_3 \mathbf{u}_3$,
and because the $\mathbf{u}_i$ are orthonormal,
$\norm{A\mathbf{x}}^2 = c_1^2 \sigma_1^2 + c_2^2 \sigma_2^2 + c_3^2 \sigma_3^2$.
This is a weighted average of $\sigma_i^2$ with weights summing to $1$, so
it is maximized when all the weight is on the largest $\sigma_i^2$ ---
i.e., $c_1 = \pm 1$ --- giving $\norm{A\mathbf{x}}^2 = \sigma_1^2 = 36$,
hence $\norm{A\mathbf{x}} = 6$.

\textbf{Trick answer:}
$\sqrt{6^2 + 3^2 + 2^2} = 7$ is the Frobenius norm $\norm{A}_F$, which
aggregates stretching across \emph{all} singular directions at once ---
it is a global measure of $A$'s ``size,'' not the maximum stretch of any
single unit input. The unit input that lands on the longest semi-axis only
gets stretched by $\sigma_1 = 6$, not by all three singular values
simultaneously.

\textbf{Trick answer:}
$36 = \sigma_1^2$ is the largest of the $\sigma_i^2$ values, which appear
as the squared semi-axis lengths in the formula
$\norm{A\mathbf{x}}^2 = \sum c_i^2 \sigma_i^2$. The maximum of
$\norm{A\mathbf{x}}^2$ is $\sigma_1^2 = 36$, but the question asks for
the maximum of $\norm{A\mathbf{x}}$, which is $\sigma_1 = 6$. A student
who picks this has confused the squared maximum with the maximum.

\textbf{Trick answer:}
$11$ just adds the singular values, $6 + 3 + 2 = 11$ --- treating them as
if they combined additively for a single input vector. They don't: the
three semi-axes lie along orthogonal directions $\mathbf{u}_1, \mathbf{u}_2, \mathbf{u}_3$,
so a single output vector can only ``use'' one of them at full strength
at a time.

\textbf{Trick answer:}
``Cannot be determined'' --- the orthogonal matrices $U$ and $V$ are needed
to find \emph{which} unit vector $\mathbf{v}_1$ achieves the maximum, but
they are not needed to find the maximum \emph{value}. The semi-axis lengths
are exactly the $\sigma_i$, and those are given.

\divider

%----------------------------------------------------------
\subsection*{Problem: Right vs left singular vectors}
%----------------------------------------------------------

A matrix $A \in \mathbb{R}^{m \times n}$ has SVD $A = U\Sigma V^T$. Which
statement most fundamentally distinguishes a right singular vector
$\mathbf{v}_i$ from a left singular vector $\mathbf{u}_i$?

\begin{enumerate}[(A)]
\item \textcolor{blue}{\textbf{$\mathbf{v}_i \in \mathbb{R}^n$ lives in the domain of $A$, while $\mathbf{u}_i \in \mathbb{R}^m$ lives in the codomain, and they are linked by $A\mathbf{v}_i = \sigma_i \mathbf{u}_i$.}}
\item $\mathbf{v}_i$ is an eigenvector of $A^TA$, while $\mathbf{u}_i$ is an eigenvector of $AA^T$.
\item $\mathbf{v}_i$ is mapped to $\sigma_i \mathbf{u}_i$ by $A$, while $\mathbf{u}_i$ is mapped to $\sigma_i \mathbf{v}_i$ by $A$.
\item $\mathbf{v}_i$ identifies a principal input direction that gets stretched by $A$, while $\mathbf{u}_i$ identifies the corresponding output semi-axis of the image ellipsoid.
\item $\mathbf{v}_i$ is the column of $V$, while $\mathbf{u}_i$ is the column of $U$.
\end{enumerate}

\textbf{Explanation:}
The most fundamental distinction between $\mathbf{v}_i$ and $\mathbf{u}_i$
is that they live in \emph{different vector spaces}:
$\mathbf{v}_i \in \mathbb{R}^n$ is in the domain, and
$\mathbf{u}_i \in \mathbb{R}^m$ is in the codomain. The relation
$A\mathbf{v}_i = \sigma_i \mathbf{u}_i$ is precisely the statement that
$A$ sends the $i$-th right singular vector to a scaled copy of the $i$-th
left singular vector. Everything else (eigenvector characterizations,
geometric stretching interpretation, columns of $U$ and $V$) is a
\emph{consequence} of this domain/codomain structure.

Geometrically: $\{\mathbf{v}_i\}$ is the orthonormal frame of input
directions on the unit sphere in $\mathbb{R}^n$; $A$ scales each by
$\sigma_i$ and sends it to $\mathbf{u}_i$, the corresponding semi-axis of
the output ellipsoid in $\mathbb{R}^m$. Whenever $m \neq n$, no
$\mathbf{v}_i$ can equal any $\mathbf{u}_j$ --- they live in spaces of
different dimension.

\textbf{Partial credit:}
``$\mathbf{v}_i$ eigenvector of $A^TA$, $\mathbf{u}_i$ eigenvector of $AA^T$''
is a true equivalent characterization, but it describes \emph{how to
compute} the singular vectors via the symmetric matrices $A^TA$ and $AA^T$
--- a derivative fact, not the fundamental distinction. A student who
picked this has internalized the eigenvalue-problem viewpoint but missed
the domain/codomain picture.

\textbf{Partial credit:}
``Principal input direction / output semi-axis'' is geometrically
correct-but-derivative --- captures the picture of input stretching to
output semi-axis, but doesn't explicitly name that the two vectors live
in \emph{different} spaces.

\textbf{Trick answer:}
``$\mathbf{v}_i$ is the column of $V$, $\mathbf{u}_i$ is the column of $U$''
is a true tautology that names \emph{where} the vectors sit in the matrices
$V$ and $U$, but says nothing about what distinguishes them mathematically.
It is the choice that ``looks right'' to a student pattern-matching on the
formula $A = U\Sigma V^T$.

\textbf{Trick answer:}
``$A$ maps $\mathbf{u}_i$ to $\sigma_i \mathbf{v}_i$'' swaps the asymmetric
mapping. The correct relation is $A\mathbf{v}_i = \sigma_i \mathbf{u}_i$
\emph{and} $A^T\mathbf{u}_i = \sigma_i \mathbf{v}_i$; the operator that
sends $\mathbf{u}_i$ back to $\sigma_i \mathbf{v}_i$ is $A^T$, not $A$.
Students who pattern-match on symmetry without checking domain/codomain
fall here.

\divider

%----------------------------------------------------------
\subsection*{Problem: Image of the unit sphere under a truncated SVD}
%----------------------------------------------------------

The matrix $A \in \mathbb{R}^{500 \times 200}$ has SVD $A = U\Sigma V^T$ with
exactly $100$ nonzero singular values
$\sigma_1 \geq \sigma_2 \geq \cdots \geq \sigma_{100} > 0$ and
$\sigma_{101} = \cdots = \sigma_{200} = 0$. Let $A_{50}$ denote the best
rank-$50$ approximation to $A$. Which of the following best describes the
image of the unit sphere in the domain under $A_{50}$?

\begin{enumerate}[(A)]
\item A solid $50$-dimensional ellipsoid in $\mathrm{span}(\mathbf{v}_{51}, \ldots, \mathbf{v}_{100}) \subset \mathbb{R}^{500}$, with semi-axis lengths $\sigma_{51}, \ldots, \sigma_{100}$.
\item A solid $50$-dimensional ellipsoid in $\mathrm{span}(\mathbf{v}_1, \ldots, \mathbf{v}_{50}) \subset \mathbb{R}^{200}$, with semi-axis lengths $\sigma_1, \ldots, \sigma_{50}$.
\item A solid $50$-dimensional ellipsoid in $\mathrm{span}(\mathbf{u}_1, \ldots, \mathbf{u}_{50}) \subset \mathbb{R}^{500}$, with semi-axis lengths $\sigma_1^2, \ldots, \sigma_{50}^2$.
\item \textcolor{blue}{\textbf{A solid $50$-dimensional ellipsoid in $\mathrm{span}(\mathbf{u}_1, \ldots, \mathbf{u}_{50}) \subset \mathbb{R}^{500}$, with semi-axis lengths $\sigma_1, \ldots, \sigma_{50}$.}}
\item A solid $50$-dimensional ellipsoid in $\mathrm{span}(\mathbf{v}_1, \ldots, \mathbf{v}_{50}) \subset \mathbb{R}^{200}$, with semi-axis lengths $\sigma_1^2, \ldots, \sigma_{50}^2$.
\end{enumerate}

\textbf{Explanation:}
The truncated SVD is
$A_{50} = \sum_{i=1}^{50} \sigma_i \mathbf{u}_i \mathbf{v}_i^T$, keeping
the $50$ largest singular values. Its image is a solid $50$-dimensional
subspace of the codomain --- namely $\mathrm{span}(\mathbf{u}_1, \ldots, \mathbf{u}_{50}) \subset \mathbb{R}^{500}$.
On the unit sphere in the domain, the action of $A_{50}$ stretches the
input direction $\mathbf{v}_i$ (for $i \leq 50$) by factor $\sigma_i$ and
sends it to $\sigma_i \mathbf{u}_i$; all directions in
$\ker(A_{50}) = \mathrm{span}(\mathbf{v}_{51}, \ldots, \mathbf{v}_{200})$
are collapsed to zero. The image is a genuine $50$-dimensional ellipsoid
with semi-axes $\sigma_1, \ldots, \sigma_{50}$.

Geometrically: the full $A$ produces a $100$-dim image ellipsoid in
$\mathbb{R}^{500}$. Truncation to rank $50$ \emph{further collapses} that
ellipsoid by zeroing the $50$ shortest \emph{nonzero} semi-axes, leaving
a solid $50$-dim ellipsoid spanned by the longest $50$ semi-axes. The $100$
already-zero singular values play no role in the truncation --- they were
already part of the kernel of $A$.

\textbf{Trick answer:}
The choice with $\mathrm{span}(\mathbf{v}_1, \ldots, \mathbf{v}_{50})$ in
$\mathbb{R}^{200}$ with semi-axes $\sigma_1^2, \ldots, \sigma_{50}^2$ is a
compound error: swaps the ambient space (domain instead of codomain) AND
uses $\sigma_i^2$ instead of $\sigma_i$. A student making both errors at
once lands here. Notably, $\sigma_i^2$ are the eigenvalues of $A^TA$, whose
eigenvectors \emph{do} live in the domain alongside the $\mathbf{v}_i$ ---
so the two errors here are mutually reinforcing rather than independent.

\textbf{Trick answer:}
The choice with $\mathrm{span}(\mathbf{v}_1, \ldots, \mathbf{v}_{50})$ in
$\mathbb{R}^{200}$ with semi-axes $\sigma_1, \ldots, \sigma_{50}$ swaps
$\mathbf{u}$ and $\mathbf{v}$. Right singular vectors $\mathbf{v}_i$ live
in the domain $\mathbb{R}^{200}$, so this answer locates the image in the
wrong space entirely. Classic domain-codomain confusion.

\textbf{Trick answer:}
The choice with semi-axes $\sigma_1^2, \ldots, \sigma_{50}^2$ in the
correct ambient space confuses singular values with eigenvalues of $A^TA$.
The semi-axis lengths are $\sigma_i$ themselves, not
$\sigma_i^2 = \lambda_i(A^TA)$. Right ambient space, wrong scaling.

\textbf{Trick answer:}
The choice with $\mathrm{span}(\mathbf{v}_{51}, \ldots, \mathbf{v}_{100})$
keeps the wrong half of the nonzero singular values --- the $50$ smallest
nonzero ones rather than the $50$ largest. ``Best rank-$50$'' keeps the
top $50$, not the next-best $50$.

\divider

%----------------------------------------------------------
\subsection*{Problem: SVD block diagram --- locating the coimage}
%----------------------------------------------------------

A diagram is given as an abstract representation of the SVD of a linear
transformation $M$, with five regions labeled (A) through (E): $U$ split
into its first $r$ columns ((A), shaded) and remaining columns (B);
$\Sigma$ as the rectangular middle block (C); $V^T$ split into its first
$r$ rows ((D), shaded) and remaining rows (E), where $r = \mathrm{rank}(M)$.
Which region best represents the \emph{coimage} of the transformation?

\begin{enumerate}[(A)]
\item Region A: shaded first $r$ columns of $U$.
\item Region B: unshaded last $m - r$ columns of $U$.
\item Region C: the diagonal $\Sigma$ block.
\item \textcolor{blue}{\textbf{Region D: shaded first $r$ rows of $V^T$.}}
\item Region E: unshaded last $n - r$ rows of $V^T$.
\end{enumerate}

\textbf{Explanation:}
The coimage $\mathrm{coim}(M) = \mathbb{R}^n / \ker(M)$ is naturally
represented in the domain by $\ker(M)^\perp$, which is spanned by the
first $r$ right singular vectors $\mathbf{v}_1, \ldots, \mathbf{v}_r$.
The rows of $V^T$ are exactly the $\mathbf{v}_i^T$, so the first $r$ rows
of $V^T$ --- region D --- house the right singular vectors that span the
coimage representative.

FTLA spine: $\mathrm{coim}(M) \cong \mathrm{im}(M)$, so regions D (in the
domain) and A (in the codomain) describe naturally isomorphic spaces. But
the \emph{coimage itself lives in the domain}, where $V$ acts --- hence D,
not A.

\textbf{Partial credit:}
Region A recognizes the rank-$r$ shaded structure and the FTLA isomorphism
$\mathrm{coim} \cong \mathrm{im}$, but places coimage on the codomain side.
A student who picked A understands the FTLA punchline but has not separated
the natural location of $\mathrm{coim}(M)$ from that of $\mathrm{im}(M)$.

\textbf{Trick answer:}
Region E is the last $n - r$ rows of $V^T$, spanning $\ker(M)$ --- the
orthogonal complement of the coimage representative. Classic
kernel-vs-coimage swap.

\textbf{Trick answer:}
Region B (unshaded columns of $U$) spans $\ker(M^T) = \mathrm{im}(M)^\perp$,
the natural representative of the \emph{cokernel}. Confuses coimage
(domain side) with cokernel (codomain side).

\textbf{Trick answer:}
Region C (the $\Sigma$ block) holds the singular values --- stretching
factors, not a subspace. A student who picks $C$ has read ``coimage'' as
``somewhere on the diagonal'' rather than as a subspace.

\divider

%----------------------------------------------------------
\subsection*{Problem: Polar decomposition --- geometric content}
%----------------------------------------------------------

A square matrix $A \in \mathbb{R}^{n \times n}$ admits a polar factorization
$A = QH$, where $Q$ is orthogonal and $H$ is symmetric positive semidefinite.
Which statement best describes the geometric content of this decomposition?

\begin{enumerate}[(A)]
\item $Q$ rotates $\mathbb{R}^n$ into the principal-axis frame, and $H$ stretches each coordinate by the corresponding singular value.
\item $H = A^T A$ and $Q = A H^{-1}$, with $H$ symmetric positive semidefinite and $Q$ orthogonal.
\item $H$ stretches by the eigenvalues of $A$ along the eigenvectors of $A$, and $Q$ realigns the result.
\item $H$ is the symmetric part $(A + A^T)/2$ and $Q$ is the orthogonal part $(A - A^T)/2$ of $A$.
\item \textcolor{blue}{\textbf{$H$ stretches along $n$ mutually orthogonal directions by the singular values of $A$, and $Q$ then rotates the stretched result.}}
\end{enumerate}

\textbf{Explanation:}
Polar decomposition writes the action of $A$ as \emph{stretch first, rotate
second}. The symmetric PSD factor $H = \sqrt{A^T A}$ is orthogonally
diagonalizable; its eigenvalues are exactly the singular values $\sigma_i$
of $A$, and its eigenvectors are the right singular vectors $\mathbf{v}_i$.
So $H$ acts on $\mathbb{R}^n$ by stretching along $n$ mutually orthogonal
axes (the $\mathbf{v}_i$) by factors $\sigma_i$. Then $Q$ rigidly rotates
the result into the codomain. This is the \textbf{fundamental} description:
$A$ is a stretch composed with a rotation, with singular values controlling
the stretch.

Cross-check via SVD: $A = U \Sigma V^T = (UV^T)(V \Sigma V^T) = QH$, with
$Q = UV^T$ and $H = V \Sigma V^T$.

\textbf{Partial credit:}
The choice ``$Q$ rotates first, $H$ stretches in the principal-axis frame''
has the right structural intuition --- rotate then stretch --- but reverses
the order; this is the \emph{left} polar decomposition $A = H'Q$, a
different factorization with a different $H'$.

\textbf{Trick answer:}
``$H$ stretches by eigenvalues of $A$ along eigenvectors of $A$'' confuses
singular values with eigenvalues and right singular vectors with
eigenvectors; $H$'s eigenvalues are $\sigma_i(A)$, not $\lambda_i(A)$, and
these agree only when $A$ is symmetric PSD.

\textbf{Trick answer:}
``$H = A^T A$, $Q = A H^{-1}$'' gets $H = \sqrt{A^T A}$ wrong as $H = A^T A$.
The matrix $A^T A$ has eigenvalues $\sigma_i^2$, not $\sigma_i$, so this
$H$ stretches by the \emph{squares} of the singular values --- the wrong
square root. Equivalently, $A H^{-1} = A (A^T A)^{-1}$ is not orthogonal
in general; the formula for $Q$ inherits the same error.

\textbf{Trick answer:}
``Symmetric/skew-symmetric decomposition'' imports the
$(A + A^T)/2 + (A - A^T)/2$ split, which has nothing to do with polar form;
the polar decomposition is multiplicative, not additive, and $(A - A^T)/2$
is skew-symmetric, not orthogonal.

\divider

\section*{WEEKS 10--11: SVD AND PCA SYNTHESIS}

\subsection*{Problem: PCA via SVD --- relationship between principal components and singular vectors}

Let $X \in \mathbb{R}^{n \times d}$ be a centered data matrix (each column has
mean zero), with SVD $X = U \Sigma V^T$. The covariance matrix is
$C = \frac{1}{n} X^T X$. Which of the following correctly describes the
relationship between the SVD of $X$ and the principal components of the
data?

\begin{enumerate}[(A)]
    \item \textcolor{blue}{\textbf{The columns of $V$ are the principal components, and the eigenvalues of $C$ are $\sigma_i^2 / n$.}}
    \item The columns of $U$ are the principal components, and the eigenvalues of $C$ are $\sigma_i^2 / n$.
    \item The columns of $V$ are the principal components, and the eigenvalues of $C$ are $\sigma_i$.
    \item The columns of $V$ are the principal components, and the eigenvalues of $C$ are $\sigma_i^2$.
    \item The columns of $U$ are the principal components, and the eigenvalues of $C$ are $\sigma_i^2$.
\end{enumerate}

\textbf{Explanation:}
This is the fundamental bridge between Week 10 (SVD) and Week 11 (PCA).
Computing directly:
$$
C = \frac{1}{n} X^T X = \frac{1}{n} (U \Sigma V^T)^T (U \Sigma V^T)
= \frac{1}{n} V \Sigma^T U^T U \Sigma V^T
= \frac{1}{n} V \Sigma^2 V^T.
$$
Reading off the eigendecomposition of $C$: the eigenvectors are the columns
of $V$ (the right singular vectors of $X$) --- these are the principal
components, the directions of maximum variance in feature space. The
eigenvalues are the diagonal entries of $\frac{1}{n}\Sigma^2$, namely
$\sigma_i^2/n$, which are precisely the variances along each PC.

The principal components live in $\mathbb{R}^d$ (feature space), so they
must be columns of $V \in \mathbb{R}^{d \times d}$, not $U \in \mathbb{R}^{n \times n}$.
The columns of $U$ have a different interpretation (they are related to
the principal component \emph{scores} after normalization).

\textbf{Partial credit:}
The choice ``columns of $V$ are PCs, eigenvalues are $\sigma_i^2$'' (without
the $1/n$) gets the eigenvectors correct and gets the squaring correct
(recognizing that $C$ involves $X^T X$, which squares singular values), but
forgets the $1/n$ normalization in the definition of the sample covariance.

\textbf{Partial credit:}
The choice ``columns of $V$ are PCs, eigenvalues are $\sigma_i$'' identifies
the principal components correctly but drops both the squaring and the
$1/n$ --- this is the relationship for $X$ itself, not for $X^T X$.

\textbf{Trick answer:}
The choice ``columns of $U$ are PCs, eigenvalues are $\sigma_i^2 / n$'' gets
the eigenvalues of $C$ correct but mislocates the principal components in
sample space rather than feature space. The eigenvalues happen to match,
which makes this distractor sticky --- but the eigenvectors are wrong
because $U$ has the wrong shape ($n \times n$ rather than $d \times d$).

\textbf{Trick answer:}
The choice ``columns of $U$ are PCs, eigenvalues are $\sigma_i^2$'' swaps
$U$ and $V$ (treating the SVD of $X$ as if it were the SVD of $X^T$) and
also drops the $1/n$. This is the most thoroughly wrong option.

\divider

\section*{WEEK 11: PRINCIPAL COMPONENT ANALYSIS}

\subsection*{Problem: Definition of the sample covariance matrix}

Let $X \in \mathbb{R}^{n \times d}$ be a centered data matrix (rows are
samples, columns are features). Which of the following expressions gives
the $d \times d$ sample covariance matrix $C$?

\begin{enumerate}[(A)]
    \item $C = X X^T$
    \item $C = \frac{1}{n} X X^T$
    \item $C = X^T X$
    \item \textcolor{blue}{\textbf{$C = \frac{1}{n} X^T X$}}
    \item $C = \frac{1}{n-1} X X^T$
\end{enumerate}

\textbf{Explanation:}
With rows as samples and columns as features, the $(i,j)$ entry of
$X^T X$ is the inner product of the $i$th and $j$th feature columns:
$$
(X^T X)_{ij} = \sum_{k=1}^n X_{ki} X_{kj}.
$$
Since the data is centered ($\sum_k X_{ki} = 0$), this is exactly $n$ times
the sample covariance between features $i$ and $j$. Dividing by $n$ gives
the covariance: $C = \frac{1}{n} X^T X$. The shape is $d \times d$,
matching the number of features.

(Note: some texts use $\frac{1}{n-1}$ for the unbiased estimator, but this
course adopts the $\frac{1}{n}$ convention --- and the question specifies
the $d \times d$ shape, which rules out the $X X^T$ variants regardless
of normalization.)

\textbf{Partial credit:}
The choice $X^T X$ has the correct shape ($d \times d$) and the correct
structure (inner products between feature columns), but omits the $1/n$
normalization. This is a standard partial-credit error: structurally
correct but dimensionally off.

\textbf{Trick answer:}
The choice $\frac{1}{n} X X^T$ produces an $n \times n$ matrix --- this
is the \emph{Gram matrix of samples}, which appears in kernel methods, but
it is not the covariance. The trick exploits surface symmetry: it has the
right $1/n$ but the wrong product order.

\textbf{Trick answer:}
$X X^T$ is the same shape error without the normalization --- doubly wrong.

\textbf{Trick answer:}
$\frac{1}{n-1} X X^T$ combines the unbiased-estimator normalization (a
real convention, just not ours) with the wrong shape. This is sticky for
students who have seen statistics elsewhere; the unfamiliar coefficient
makes it look authoritative, but the shape betrays it.

\divider

\subsection*{Problem: Correlation as cosine similarity}

Let $X \in \mathbb{R}^{n \times d}$ be a centered data matrix, and let $x_i$
and $x_j$ denote the $i$th and $j$th columns of $X$ (i.e., the centered
data for features $i$ and $j$). The Pearson correlation between features
$i$ and $j$ is
$$
\rho_{ij} = \frac{x_i \cdot x_j}{\norm{x_i}\, \norm{x_j}}.
$$
Which of the following best describes the geometric meaning of $\rho_{ij}$?

\begin{enumerate}[(A)]
    \item \textcolor{blue}{\textbf{$\rho_{ij}$ is the cosine of the angle between the feature vectors $x_i$ and $x_j$ in $\mathbb{R}^n$.}}
    \item $\rho_{ij}$ is the cosine of the angle between $x_i$ and $x_j$ in $\mathbb{R}^d$.
    \item $\rho_{ij}$ is the projection of $x_i$ onto $x_j$.
    \item $\rho_{ij}$ is the distance between $x_i$ and $x_j$ in feature space.
    \item $\rho_{ij}$ is the variance of the joint distribution of features $i$ and $j$.
\end{enumerate}

\textbf{Explanation:}
The formula $\rho_{ij} = (x_i \cdot x_j) / (\norm{x_i} \norm{x_j})$ is
\emph{exactly} the cosine of the angle between $x_i$ and $x_j$ --- this is
the definition of cosine similarity from Week 5. The feature vectors
$x_i, x_j \in \mathbb{R}^n$ live in sample space (one component per sample),
so the angle is measured in $\mathbb{R}^n$.

This is the conceptual punchline of correlation: \emph{correlation is just
cosine similarity applied to centered feature vectors.} Centering makes
the inner product equal to the covariance (up to $n$); normalizing by the
norms turns covariance into correlation, which is bounded in $[-1, 1]$
exactly because cosine is.

\textbf{Trick answer:}
The choice with $\mathbb{R}^d$ instead of $\mathbb{R}^n$ swaps sample
space and feature space --- a classic D7-style type confusion. Each $x_i$
has one entry per sample, so it lives in $\mathbb{R}^n$, not $\mathbb{R}^d$.

\textbf{Trick answer:}
``Projection of $x_i$ onto $x_j$'' confuses correlation with the projection
coefficient $(x_i \cdot x_j)/\norm{x_j}^2$. The projection coefficient is
not symmetric in $i$ and $j$, but correlation is.

\textbf{Trick answer:}
``Distance between $x_i$ and $x_j$'' confuses inner products (angles) with
norms of differences (distances). These are different geometric objects:
$\rho_{ij} = 1$ does not mean $x_i = x_j$; it means they point in the
same direction.

\textbf{Trick answer:}
``Variance of the joint distribution'' is the most thoroughly statistical
distractor --- it sounds plausible because correlation is a statistical
quantity, but variance is not bounded in $[-1, 1]$ and is a property of a
single random variable, not a relationship between two.

\divider

\section*{WEEKS 5, 6, 12 SYNTHESIS}

\subsection*{Problem: Geometric role of the perceptron weight vector}

A perceptron with weights $\mathbf{w} \in \mathbb{R}^n$, bias
$b \in \mathbb{R}$, and sigmoid activation $\sigma$ computes
$y = \sigma(\mathbf{w}^T \mathbf{x} + b)$, with decision boundary
$\{\mathbf{x} : \mathbf{w}^T \mathbf{x} + b = 0\}$. Which statement best
and most fundamentally describes the geometric role of $\mathbf{w}$?

\begin{enumerate}[(A)]
    \item $\mathbf{w}$ points from the origin toward the centroid of the training inputs whose target label is close to $1$.
    \item $\mathbf{w}$ lies in the decision boundary, since $\mathbf{w}$ and $\mathbf{x}$ are vectors in the same space $\mathbb{R}^n$.
    \item \textcolor{blue}{\textbf{$\mathbf{w}$ is normal to the decision boundary, and $\mathbf{w}^T \mathbf{x} + b$ equals $\norm{\mathbf{w}}$ times the signed distance from $\mathbf{x}$ to that boundary.}}
    \item $\mathbf{w}$ determines the decision boundary.
    \item $\mathbf{w}$ is the direction of steepest ascent of the pre-activation $\mathbf{x} \mapsto \mathbf{w}^T \mathbf{x} + b$.
\end{enumerate}

\textbf{Explanation:}
Two facts together pin down the geometric role of $\mathbf{w}$.

\emph{(i) $\mathbf{w}$ is normal to the boundary.} The decision boundary
is the affine hyperplane $\{\mathbf{x} : \mathbf{w}^T \mathbf{x} + b = 0\}$.
For any two points $\mathbf{x}_1, \mathbf{x}_2$ on this boundary,
$\mathbf{w}^T (\mathbf{x}_1 - \mathbf{x}_2) = 0$, so $\mathbf{w}$ is
orthogonal to every direction \emph{along} the boundary --- i.e.,
$\mathbf{w}$ is a normal vector to the hyperplane.

\emph{(ii) The pre-activation magnitude has metric meaning.} Decompose
any input as
$\mathbf{x} = \mathbf{x}_\parallel + t \cdot \mathbf{w}/\norm{\mathbf{w}}$,
where $\mathbf{x}_\parallel$ is the orthogonal projection of $\mathbf{x}$
onto the boundary and $t$ is the signed distance from $\mathbf{x}$ to
the boundary along $\mathbf{w}$. Then
$$
\mathbf{w}^T \mathbf{x} + b
= \mathbf{w}^T \mathbf{x}_\parallel + b + t \norm{\mathbf{w}}
= t \norm{\mathbf{w}},
$$
since $\mathbf{w}^T \mathbf{x}_\parallel + b = 0$ on the boundary. So the
pre-activation is literally $\norm{\mathbf{w}}$ times the signed distance
from $\mathbf{x}$ to the decision boundary.

This is the synthesis. The perceptron is a geometric construction in
disguise: the dot product $\mathbf{w}^T \mathbf{x}$ from Week 5, the
orthogonal-projection picture from Week 6, and the affine-hyperplane
picture all collapse together. ``Classify $\mathbf{x}$'' means ``compute
the signed distance from $\mathbf{x}$ to a hyperplane and squash through
$\sigma$.''

\textbf{Partial credit:}
``$\mathbf{w}$ determines the decision boundary'' is true but
incomplete --- it captures only that $\mathbf{w}$ specifies the
hyperplane, missing the second fact that the \emph{magnitude}
$\mathbf{w}^T \mathbf{x} + b$ also has metric meaning as a scaled
signed distance. One rung below the correct answer on the ladder of
geometric content.

\textbf{Trick answer:}
``$\mathbf{w}$ is the direction of steepest ascent of the pre-activation''
is \emph{true}, but it is a calculus consequence of the linear-algebraic
fact that level sets of $\mathbf{w}^T \mathbf{x} + b$ are hyperplanes
normal to $\mathbf{w}$. The gradient fact describes a consequence of
$\mathbf{w}$'s geometric role; it does not capture what $\mathbf{w}$
geometrically \emph{is}. A true statement that fails to be the
\emph{best and most fundamental} description.

\textbf{Trick answer:}
``$\mathbf{w}$ lies in the decision boundary, since $\mathbf{w}$ and
$\mathbf{x}$ are vectors in the same space'' is an input/weight-space
confusion. Living in the same ambient $\mathbb{R}^n$ does not put
$\mathbf{w}$ on the boundary --- $\mathbf{w}$ generically does not
satisfy $\mathbf{w}^T \mathbf{w} + b = 0$. The fact that $\mathbf{w}$ is
\emph{normal} to the boundary is precisely the opposite of lying in it.

\textbf{Trick answer:}
``$\mathbf{w}$ points toward the centroid of the positive class'' is the
folk picture: the weight vector ``points toward the positive examples.''
Directionally suggestive but mathematically false in general --- the
centroid of positive examples need not lie along $\mathbf{w}$, and after
training $\mathbf{w}$ is determined by margin geometry, not class
centroids. This story is a half-remembered description of LDA or a
class-mean classifier, not a perceptron.

\divider

\section*{WEEK 12: NEURAL NETWORKS}

\subsection*{Problem: Backpropagation as the chain rule}

In a feed-forward network, layer $k$ takes input $\mathbf{h}_{k-1}$ and
produces $\mathbf{z}_k = W_k \mathbf{h}_{k-1} + \mathbf{b}_k$ followed by
$\mathbf{h}_k = \sigma(\mathbf{z}_k)$, with $\sigma$ applied componentwise.
Define $\boldsymbol{\delta}_k = \partial L / \partial \mathbf{z}_k$.

Backpropagation produces a recursion that expresses $\boldsymbol{\delta}_k$
in terms of $\boldsymbol{\delta}_{k+1}$. Which is the most accurate
description of \emph{why} this recursion takes the form it does?

\begin{enumerate}[(A)]
\item Since the loss is computed at the output, a hidden layer's contribution decays geometrically with depth, and the recursion encodes that decay.
\item \textcolor{blue}{\textbf{The map $\mathbf{z}_k \mapsto \mathbf{z}_{k+1}$ factors through $\mathbf{h}_k$, and its derivative is the product of the derivatives of the two pieces --- a diagonal matrix from componentwise $\sigma$ and the matrix $W_{k+1}$ from the affine step.}}
\item The recursion is a consequence of gradient descent on the loss, in which the learning rate is replaced by $\mathrm{diag}(\sigma'(\mathbf{z}_k))$ to keep updates well-scaled across layers.
\item Backpropagation distributes the output error $\mathbf{h}_{\text{out}} - \mathbf{t}$ across layers in proportion to each layer's share of the network's total weight magnitude.
\item Each layer must be inverted to send information from the output back toward the input, so the recursion involves $W_{k+1}^{-1}$ acting on $\boldsymbol{\delta}_{k+1}$.
\end{enumerate}

\textbf{Explanation:}
Backpropagation is the chain rule applied to a composition. The map from
$\mathbf{z}_k$ to $\mathbf{z}_{k+1}$ factors through $\mathbf{h}_k$:
componentwise activation $\sigma$ takes $\mathbf{z}_k$ to $\mathbf{h}_k$,
and the affine map $W_{k+1} \cdot + \mathbf{b}_{k+1}$ takes $\mathbf{h}_k$
to $\mathbf{z}_{k+1}$. The chain rule operates on \emph{derivatives}, and
each piece has a well-defined linear-map derivative: the derivative of the
componentwise $\sigma$ is a diagonal matrix (because applying $\sigma$
componentwise has no cross-coordinate effect), and the derivative of the
affine map is just the matrix $W_{k+1}$ (the bias is a constant and drops out).

To pull the gradient $\boldsymbol{\delta}_{k+1} = \partial L / \partial \mathbf{z}_{k+1}$
back through this composition, we apply the chain rule. The structurally
important fact --- the linear-algebra heart of the entire course showing up
in the deep-learning setting --- is that pulling a gradient back through a
linear map multiplies by its \emph{transpose}, not its inverse. (This is
the same transpose that drives least squares, the FTLA, and the
pseudoinverse: forward through a linear map uses $W$; backward through it
uses $W^T$.) Composing the transposes of the two pieces in reverse order
recovers $\boldsymbol{\delta}_k$ in terms of $\boldsymbol{\delta}_{k+1}$,
with the diagonal $\sigma'$ matrix and the matrix transpose $W_{k+1}^T$
both appearing. This is what the correct answer states in words.

\textbf{Trick answer:}
``Geometric decay with depth'' is a possible \emph{training-dynamics}
phenomenon (the vanishing-gradient effect), not a property of the
recursion itself. The recursion can shrink, preserve, or amplify
$\boldsymbol{\delta}$ depending on the weight matrix and the activation
slopes; geometric decay is one specific regime, not the structural form.

\textbf{Trick answer:}
``Distributes the output error in proportion to weight magnitude'' is a
folk picture of error ``flowing back'' --- no chain rule, and the
proportional-to-weight-magnitude story is not what the recursion does.
This is the kind of plausible-sounding pseudo-explanation that the chain
rule replaces.

\textbf{Trick answer:}
``Replaces the learning rate with $\mathrm{diag}(\sigma'(\mathbf{z}_k))$''
confuses the recursion with the gradient-descent update step. The gradient
$\boldsymbol{\delta}_k$ is computed before any learning rate enters; the
learning rate appears later, when $\boldsymbol{\delta}_k$ is used to form
$\partial L / \partial W_k$ and update weights.

\textbf{Trick answer:}
``Each layer must be inverted, so the recursion involves $W_{k+1}^{-1}$'' is
the naive ``to go backward, invert'' instinct. Wrong on the linear
algebra: weight matrices are generally not square, so $W_{k+1}^{-1}$ is
not even defined; and even when they are square, the chain rule produces
the transpose, not the inverse. Adjoint and inverse coincide only for
orthogonal matrices --- a very special case.

\divider

\subsection*{Problem: Why nonlinearity is required}

A feed-forward neural network alternates affine layers
$\mathbf{z}_k = W_k \mathbf{h}_{k-1} + \mathbf{b}_k$ with a componentwise
nonlinear activation $\mathbf{h}_k = \sigma(\mathbf{z}_k)$. Removing the
activations from such a network --- so each layer is purely affine and
the layers are composed directly --- collapses the network's expressive
power dramatically. Which of the following best describes the
\emph{fundamental} reason a nonlinear $\sigma$ is required?

\begin{enumerate}[(A)]
\item Without $\sigma$, the Universal Approximation Theorem does not apply, so the network is not guaranteed to approximate arbitrary target functions.
\item Without $\sigma$, gradients cannot flow back through the network during backpropagation, so the weights cannot be trained.
\item Without $\sigma$, hidden activations are unbounded and the network's outputs cannot be interpreted as probabilities or classifications.
\item Without $\sigma$, the network can only fit linear targets, and most real-world relationships between inputs and outputs are nonlinear.
\item \textcolor{blue}{\textbf{Without $\sigma$, the entire network --- regardless of depth or width --- represents only an affine function of its input, so no amount of stacking enlarges the class of functions the network can compute.}}
\end{enumerate}

\textbf{Explanation:}
The structural fact is that the composition of affine maps is affine.
Writing two consecutive layers without activation,
$$\mathbf{z}_2 = W_2(W_1 \mathbf{x} + \mathbf{b}_1) + \mathbf{b}_2 = (W_2 W_1) \mathbf{x} + (W_2 \mathbf{b}_1 + \mathbf{b}_2),$$
the result is again of the form $W' \mathbf{x} + \mathbf{b}'$. Iterating,
a network of any depth without $\sigma$ collapses to a single affine map.
Depth gives no new expressive power; width gives no new expressive power;
the function class is exactly the affine maps. The nonlinear $\sigma$ is
the \emph{only} ingredient that breaks this collapse and makes depth
meaningful. Every other consequence --- universal approximation, the
ability to fit nonlinear targets, the bounded-output interpretation of
sigmoid specifically --- is downstream of this structural fact.

This is the foundational ``why $\sigma$ exists'' answer for the course.
The other options each name something true \emph{about} networks with
$\sigma$, but they describe consequences rather than the structural cause.

\textbf{Trick answer:}
``Gradients cannot flow back through the network'' imports a backpropagation
concern. Affine maps are perfectly differentiable, and backpropagation
through a network with no activations works mechanically: it just computes
the gradient of an affine map, which exists and is well-defined. The
training pathology if you removed $\sigma$ is not that gradients vanish
but that the only function the network \emph{can} represent is affine.

\textbf{Trick answer:}
``Hidden activations are unbounded'' confuses the role of nonlinearity in
general with the specific behavior of sigmoid. Sigmoid bounds outputs to
$(0,1)$; tanh bounds to $(-1,1)$; ReLU does not bound at all. Boundedness
is a feature of \emph{particular} activations, not the structural reason
any nonlinearity is needed. A network using ReLU still escapes the
affine-collapse problem despite ReLU being unbounded above.

\textbf{Trick answer:}
``UAT does not apply'' restates the question in different vocabulary. The
Universal Approximation Theorem \emph{assumes} a nonlinear activation as
a hypothesis, so saying ``without $\sigma$ UAT fails'' is just naming
UAT's hypothesis, not explaining why $\sigma$ is needed. The genuine
answer addresses the structural cause that makes UAT's hypothesis necessary
in the first place.

\textbf{Trick answer:}
``Can only fit linear targets'' locates the problem in the world rather
than in the network. Even if every target function in the universe were
linear, removing $\sigma$ would still mean the network's depth and width
are wasted: a 50-layer linear regressor would do exactly the same thing
as a 1-layer linear regressor. The issue is intrinsic to the network
architecture, not contingent on the target.

\divider

\end{document}
